% =============================================================================
% Do Mortality Prediction Intervals Survive a Pandemic?
% A Pre-Registered Calibration Audit of Ten Forecasting Families and Seven
% Uncertainty Mechanisms Across the COVID-19 Break
%
% Build with `make` (latexmk -pdf) from paper/.
%
% CLASS NOTE. Target venues are ASTIN Bulletin and Annals of Actuarial Science
% (Cambridge University Press). Cambridge's own journal classes are not freely
% redistributable, so this draft uses the standard `article' class at 11pt with
% natbib author-year citations, which is what CUP accepts for submission in any
% case. Swapping to the CUP template at submission is a preamble-only change.
% =============================================================================
\documentclass[11pt,a4paper]{article}

% --- engine-conditional font setup -------------------------------------------
% pdfTeX (latexmk -pdf, the intended route) gets T1/utf8/lmodern/microtype.
% XeTeX/LuaTeX (the tectonic fallback in the Makefile) skip the pdfTeX-only
% packages. Sources are kept ASCII-only so both routes render identically.
\usepackage{iftex}
\ifPDFTeX
  \usepackage[T1]{fontenc}
  \usepackage[utf8]{inputenc}
  \usepackage{lmodern}
  \usepackage{microtype}
\fi

% A file path in \texttt{} is one unbreakable box; a justified line carrying
% one overhangs the margin (3-20 pt, measured on the XeTeX route, where the
% pdfTeX-only microtype above is unavailable). Let TeX stretch the line
% instead in its final pass -- engine-independent, unlike microtype.
\setlength{\emergencystretch}{3em}

\usepackage[a4paper,margin=2.5cm,marginparwidth=2.0cm,marginparsep=0.35cm]{geometry}
\usepackage{amsmath,amssymb,amsthm}
\usepackage{booktabs}
\usepackage{graphicx}
\usepackage{xcolor}
\usepackage{caption}
\usepackage{threeparttable}
\usepackage{longtable}
\usepackage{array}
\usepackage{enumitem}
\usepackage[authoryear,round]{natbib}
\usepackage[colorlinks=true,linkcolor=blue!50!black,citecolor=blue!50!black,urlcolor=blue!50!black]{hyperref}

\graphicspath{{figures/}}
\setlist{nosep}
\captionsetup{font=small,labelfont=bf,width=0.92\textwidth}

% --- \todo: red inline marker + red margin note, numbered ----------------------
% \todo{...}       body text only (uses \marginpar; illegal inside floats,
%                  captions, footnotes, headings)
% \todoinline{...} same marker without the margin note; safe anywhere
\newcounter{todocount}
\newcommand{\todoinline}[1]{%
  \stepcounter{todocount}%
  \textcolor{red}{\textbf{[TODO~\thetodocount: #1]}}}
\newcommand{\todo}[1]{%
  \todoinline{#1}%
  \marginpar{\color{red}\scriptsize\raggedright\textbf{TODO~\thetodocount.} #1}}

% --- \inputtable{name}: pull a generated fragment, or a placeholder --------------
% Tables are written by scripts/make_tables.py (see tables/README.md). A
% missing fragment renders as a placeholder so the draft always compiles.
\newcommand{\inputtable}[1]{%
  \IfFileExists{tables/#1.tex}%
    {\input{tables/#1}}%
    {\todoinline{table \texttt{#1} is produced by \texttt{scripts/make\_tables.py} from \texttt{results/*.parquet}; not yet generated}}}

% --- \inputfigure{name}: \includegraphics of a generated figure, or a placeholder
% Figures are written by the analysis stage as figures/<name>.pdf. A missing
% file renders as a placeholder so the draft always compiles. Optional
% argument = width (default \textwidth).
\newcommand{\inputfigure}[2][\textwidth]{%
  \IfFileExists{figures/#2.pdf}%
    {\includegraphics[width=#1,height=0.88\textheight,keepaspectratio]{figures/#2.pdf}}%
    {\todoinline{figure \texttt{#2.pdf} is produced by the analysis stage from \texttt{results/*.parquet}; not yet generated}}}

% --- notation ------------------------------------------------------------------
\newcommand{\E}{\mathbb{E}}
\newcommand{\Prob}{\mathbb{P}}
\newcommand{\ind}{\mathbb{1}}
\newcommand{\CRPS}{\operatorname{CRPS}}
\newcommand{\LS}{\operatorname{LS}}
\newcommand{\IS}{\operatorname{IS}}
\newcommand{\RMSE}{\operatorname{RMSE}}
\newcommand{\MAE}{\operatorname{MAE}}
\newcommand{\PIT}{\operatorname{PIT}}
\newcommand{\annuity}{\ddot{a}_{65}}
\newcommand{\Hone}{\textbf{H1}}
\newcommand{\Htwo}{\textbf{H2}}
\newcommand{\Hthree}{\textbf{H3}}
\newcommand{\Hfour}{\textbf{H4}}
\newcommand{\Hfive}{\textbf{H5}}
\newcommand{\pkg}[1]{\texttt{#1}}

\theoremstyle{definition}
\newtheorem{hypothesis}{Hypothesis}

\title{Do Mortality Prediction Intervals Survive a Pandemic?\\[0.4em]
  \large A Pre-Registered Calibration Audit of Ten Forecasting Families and
  Seven Uncertainty Mechanisms Across the COVID-19 Break}
\author{%
  Hamid Jahani\thanks{Corresponding author.
    Email address: \texttt{hamid.jahani@modares.ac.ir} (Hamid Jahani).}\\[0.45em]
  \normalsize Department of Data Science, Faculty of Interdisciplinary Science and Technology,\\
  \normalsize Tarbiat Modares University, Tehran, Iran}
\date{5 September 2026}

\begin{document}
\maketitle

\begin{abstract}
Annuity and pension valuations price the \emph{spread} of a mortality
forecast, not only its centre: reserves, risk margins and longevity capital
all rest on a nominal 95\% interval covering the realised outcome about 95\%
of the time. Such intervals have been produced since \citet{lee1992modeling},
and lately by neural extensions of Lee--Carter, but never audited against a
structural break paired with a quiet control that separates break-induced
failure from standing miscalibration. We pre-register that audit on the June
2026 Human Mortality Database vintage, the first with final data through
2024. Ten forecasting families crossed with seven uncertainty
mechanisms---model-native, Poisson bootstrap, deep ensembles, Monte Carlo
dropout, split conformal, EnbPI and a joint-path copula band---give fifty
admissible cells, all scored on one code path, across three regimes: a quiet
expanding-origin control (1990--2014), the COVID break (train to 2019, test
2020--2024) and a 1914--1922 placebo.

The control carries the finding: most miscalibration predates the pandemic.
Native intervals from the Lee--Carter lineage already cover far below
nominal in quiet decades---Lee--Carter 0.702, Poisson Lee--Carter 0.682,
Renshaw--Haberman 0.731 against 0.950---and the registered one-sided
hypothesis of universal under-coverage across the break is refuted: its
two-sided restatement grows in only 27 of 50 cells, because the break mainly
sorts arms by the side of nominal they began on, over-covering ones moving
\emph{towards} it. Benchmarked against independence at each arm's own
marginal rate, joint path coverage is \emph{higher} in 50 of 50 control and
50 of 50 break cells: horizons are strongly positively dependent, a
dependence noted in the literature but never measured. The rate-level
figures understate what happens in money---in the quiet control the native
Lee--Carter interval covers 0.253 of realised annuity factors and the LSTM
deep ensemble 0.157. Most of these intervals were never calibrated to begin
with, and the pandemic mainly re-orders them.
\end{abstract}

\noindent\textbf{Keywords:} mortality forecasting; longevity risk;
calibration; prediction intervals; proper scoring rules; conformal
prediction; COVID-19; Lee--Carter; neural networks; pre-registration.

\vspace{0.5em}
\noindent\textbf{Pre-registration.} \texttt{PREREGISTRATION.md}, committed
2026-08-25 before any model was fit to real data, with four registered
addenda, \texttt{PREREGISTRATION-ADDENDUM-1..4.md}: addendum~1
(data-prerequisite strata and sensitivities), 2 (scoring-target
clarifications) and 3 (defect handling, scoring repairs and hypothesis
status, including the two-sided restatement of \Htwo{} and the re-specified
\Hthree{} comparator of \S8), all committed 2026-08-26 and still before any
real-data fit, and 4 (a sixty-year training-window cap for the Gaussian
process), committed 2026-08-29 during the sweeps and before any GP cell on a
longer panel had produced a result. SHA-256, in that order:\\
\texttt{\small bbbe3860a5446d063e186e56555afa893b07c36cf15a990d07567471affe50be}\\
\texttt{\small 8b43083afae9f0eaafce5e89381349d74cbbafb94022b56d1f2f39032a308889}\\
\texttt{\small 1618e9fdcc1d6cbc760ce40ea496c1285fe44b8fc7520e8687e2101343ebdbc1}\\
\texttt{\small cca2881a3a7ae6a17bbec7a4c6440192a6f9aa1cdd5766c13ef0aaec8ddeb16f}\\
\texttt{\small 586e01780cceeea32878d1a493662c2883d017396459dc3442a6053f101a564b}

\section{Introduction}
\label{sec:intro}

\subsection{Longevity risk is priced off the interval}

A life annuity is a bet on a distribution. The best-estimate reserve of an
annuity book is a functional of the central mortality projection, but
everything that sits on top of it---the risk margin, the longevity component
of solvency capital, the price at which a longevity swap or a bulk annuity
transfer clears---is a functional of the projection's \emph{spread}. Under
the Solvency~II standard formula the longevity stress is a permanent
20\% fall in mortality rates \citep{eiopa2015delegated}; under an internal
model it is the 99.5\% one-year quantile of the change in own funds, which a
stochastic mortality model must supply from its own predictive distribution.
Either way the quantity carried on the balance sheet is a tail of a forecast
distribution, and its adequacy rests on one property: that a nominal 95\%
interval actually contains the realised outcome about 95\% of the time. An
interval that is too narrow under-reserves; one that is too wide over-prices
the hedge. Neither error is visible in a point-accuracy comparison.

Stochastic mortality models have produced intervals from the start.
\citet{lee1992modeling} forecast the period index $\kappa_t$ by a random walk
with drift and reported prediction intervals for life expectancy; the
Cairns--Blake--Dowd programme built a family of models expressly to generate
density forecasts \citep{cairns2006two,cairns2009quantitative,cairns2011mortality};
the actuarial toolkit \pkg{StMoMo} \citep{villegas2018stmomo} exposes bootstrap
and simulation intervals for all of them. The neural extensions of
Lee--Carter that followed
\citep{richman2021neural,nigri2019deep,perla2021time,scognamiglio2022calibrating}
improved point accuracy across many populations, and several now emit
intervals of their own
\citep{marino2023neural,miyata2022extending,schnurch2022point}. Whether
those intervals are \emph{calibrated} is a separate question from whether
they exist, and it is the question that decides what they are worth to an
actuary.

\subsection{The break that makes the question answerable}

Calibration is testable only where the truth is eventually observed, and the
test is informative only where the truth is allowed to surprise the model.
From the mid-1950s until 2019 the mortality surfaces of the high-income
populations in the Human Mortality Database (HMD) improved smoothly enough
that almost any trend extrapolation looked calibrated in sample. The years
2020--2024 broke that surface: period life expectancy fell in 2020 and 2021
in most HMD countries \citep{aburto2022quantifying}
\todo{verify magnitude and country coverage against Aburto et al.\ (2022)
before quoting any figure}, and the subsequent recovery was uneven by country
and by age. Any prediction interval constructed from data ending in 2019 was,
for the first time in the modern record, put to a genuine out-of-sample test
by reality.

The result of that test has not been scored. The closest study,
\citet{barigou2023bayesian}, has an exemplary evaluation
design---leave-future-out validation, log score and CRPS---but its pandemic
analysis is a simulation: French male deaths perturbed by ``two years of
excess mortality followed by one year of lower mortality'', because in 2021
real post-shock data did not exist. They exist now. The HMD vintage of
15~June~2026 (Methods Protocol v6) is the first with final single-age data
through 2024 for twenty populations. This paper uses it.

\subsection{What the literature measures}

We searched the full text of the mortality-forecasting library assembled for
this study (about twenty modelling papers; the term counts and quotations are
recorded in the project repository) for what each paper actually reports.
The counts below are over that library; where wider reading changed them,
the correction is stated.

\begin{itemize}
\item \emph{Point accuracy is the default and often the only metric.}
  \citet{richman2021neural}, the multi-population neural Lee--Carter that
  later work builds on, is fit to every HMD country and evaluated on mean
  squared error alone: the term ``MSE'' appears twenty-eight times in the
  full text; ``coverage'', ``CRPS'' and ``log score'' appear zero times.
  \citet{scognamiglio2022calibrating}, \citet{nigri2019deep},
  \citet{petnehazi2019mortality} and \citet{perla2021time} have the same
  profile.
\item \emph{Intervals are drawn but not checked.} \citet{miyata2022extending}
  advertise ``the ability to forecast with confidence intervals'', plot
  97.5\% intervals for the latent factor and for mortality at ages 30--80
  (their Section~6.4 and Figures~7--8), and then compare against Lee--Carter
  by MSE only (their Table~3). The paper never asks whether observed
  mortality falls inside the intervals it draws. The same authors note, of
  \citet{schnurch2022point}, that the intervals there are ``not based on the
  endogenous randomness \ldots\ but on the exogenous randomness driven by the
  random seed for the NN'': variance across seeds, which measures nothing
  about predictive coverage.
  \todo{Schn\"urch and Korn (2022) not yet read; confirm this characterisation
  and whether they tabulate coverage}
\item \emph{Where coverage is reported, it is reported without a sharpness
  penalty.} \citet{marino2023neural} is, to our knowledge, the one neural
  Lee--Carter paper that computes interval coverage. It reports the
  prediction-interval coverage probability (PICP) and the mean interval width
  (MPIW) as two separate indicators and treats $\mathrm{PICP}=1$ as the best
  outcome: the LSTM ``offers always a greater probability coverage, in most
  cases due to the PI width''. An interval wide enough to contain everything
  scores~$1$; without a width-penalised criterion, ``wider'' and ``better''
  are indistinguishable. Buried in the same paper, and not followed up, is
  the observation that the classical ARIMA intervals attain empirical
  coverage ``around 50\%'' and, for Spanish males, ``stable around
  33\%''---nominal 95\% Lee--Carter intervals covering a third to a half of
  realisations. The test window (2001--2018, three countries) contains no
  structural break. The most recent competitor, the gradient-boosted
  multi-population model of \citet{miao2025gradient}, repeats the pattern one
  level up: marginal coverages of 90\%, 94\% and 96\% against a nominal 95\%
  are reported and the closest is declared ``the most properly calibrated'',
  again with no proper score and no width penalty alongside.
\item \emph{Ex-post density evaluation exists, but predates both neural
  models and the break.} \citet{cairns2011mortality} assess the density
  forecasts of six stochastic models \emph{ex ante}, by the biological
  reasonableness of their fan charts; there is no PIT, no coverage
  computation, no proper score and no held-out evaluation.
  \citet{dowd2010backtesting} do backtest multi-period-ahead density
  forecasts ex post, for six classical models on English and Welsh males.
  \todo{VERIFY against Dowd et al.\ (2010b) once read: which backtests, whether
  hit rates or PIT values are tabulated, horizons, ages, and the verdict}
  That work is the direct ancestor of the present audit; it is confined to
  one population, to models with native intervals, and to a period without a
  break.
\item \emph{Absent from mortality forecasting, so far as we can find:} PIT or
  rank histograms; conformal prediction; joint (path) coverage over horizons;
  and, with the single exception of \citet{shang2018model}, formal
  forecast-comparison testing (Diebold--Mariano, model confidence sets).
  \todo{Crossref/Scholar search for ``conformal'' + ``mortality'' before
  submission; the library sweep was arXiv-biased}
\end{itemize}

The principle that the literature's coverage practice violates is old and
precise: a probabilistic forecaster should \emph{maximise the sharpness of
the predictive distributions subject to calibration}
\citep{gneiting2007probabilistic}. Coverage without width is not
calibration; width without coverage is not sharpness; and neither is a
proper score \citep{gneiting2007strictly}. The audit below scores every
forecast by both, and by the proper scores that combine them.

\subsection{This study}

The question is whether the prediction intervals of stochastic,
machine-learning and neural mortality forecasting models attain their
nominal coverage when the evaluation window crosses the COVID-19 break, and
whether the answer depends on the model family, on the uncertainty
mechanism, or on both. The last clause is the design. Existing comparisons
vary the model and inherit whatever uncertainty machinery each model happens
to ship with, so a coverage failure cannot be attributed: is the
architecture wrong about the future, or is the interval construction wrong
about the architecture? We treat \emph{model family} and \emph{uncertainty
mechanism} as separate crossed factors---ten families, seven mechanisms,
fifty admissible cells---so that the same architecture is audited under
several mechanisms and the same mechanism across several architectures.
That separation is what makes this a study rather than a bake-off.

Three regimes separate ordinary forecast error from break-induced failure: a
\emph{stable} control (expanding origins 1990--2014, every test year before
2020); the \emph{shift} regime (train to 2019, test 2020--2024); and a
\emph{placebo} break (train to 1913, test 1914--1922, spanning the First
World War and the 1918 influenza pandemic) that asks whether whatever we
find is specific to COVID-19 or is simply what intervals do at any break.

Five hypotheses were fixed and hash-stamped before any model touched real
data. They are reproduced verbatim from the pre-registration.

\begin{hypothesis}[H1]
Model rankings by RMSE and by CRPS disagree (rank correlation $<1$; at least
one top-3 inversion) in the shift regime.
\end{hypothesis}
\begin{hypothesis}[H2]
Nominal 95\% intervals under-cover in the shift regime for every model
family; the shortfall exceeds the stable-regime shortfall.
\end{hypothesis}
\begin{hypothesis}[H3]
Joint path coverage over $h = 1,\dots,5$ is materially below marginal
coverage at the same nominal level for every family (reported gap; test:
paired difference $>0$).
\end{hypothesis}
\begin{hypothesis}[H4]
Coverage failure is non-uniform in age: shift-regime coverage at ages 65--99
is lower than at ages 25--64 for every family.
\end{hypothesis}
\begin{hypothesis}[H5]
Interval miscalibration propagates to derived quantities: empirical coverage
of 95\% intervals for $e_{65}$ and annuity factor $\annuity$ (2\%) departs
from nominal by more than the corresponding stable-regime departure.
\end{hypothesis}

No directional claim is registered about neural-versus-classical ordering;
the crossed design estimates it.

\subsection{Contributions}

This paper proposes no new mortality model. Its contributions are:
\begin{enumerate}
\item \emph{The audit.} The first evaluation of probabilistic mortality
  forecasts---classical, neural and conformal---against realised 2020--2024
  mortality across twenty populations, with every interval constructed
  strictly from pre-2020 data.
\item \emph{The crossed design.} Model family and uncertainty mechanism as
  separate factors, so that a coverage failure can be attributed to the
  architecture, to the mechanism, or to their interaction. Claims are stated
  as contrasts within admissible sub-grids, never as main effects over a
  ragged grid.
\item \emph{Resolution of the failure.} By horizon (marginal versus joint
  path coverage, \Hthree), by age (\Hfour) and in economic terms (life
  expectancy and annuity-factor intervals, \Hfive), with the 1914--1922
  placebo as a second event.
\item \emph{The protocol.} A pre-registered, synthetic-truth-validated,
  oracle-parity-checked evaluation harness in which every model emits
  predictive samples through one interface and is scored through one code
  path, with cluster-bootstrap forecast-comparison inference that respects
  the fact that a pandemic is one shock, not twenty. The harness is reusable
  for any future break.
\end{enumerate}

Section~\ref{sec:related} places the audit in the literature.
Section~\ref{sec:data} describes the data and the conventions that make
Poisson scoring of fractional deaths well defined. Section~\ref{sec:design}
sets out the regimes, the grid, the interface and the validation gates.
Section~\ref{sec:metrics} defines every metric. Sections~\ref{sec:results}
and~\ref{sec:actuarial} are the results and their actuarial reading; both
are placeholders in this draft. Section~\ref{sec:discussion} pre-commits the
limitations and Section~\ref{sec:conclusion} concludes.

\section{Related work}
\label{sec:related}

The literature was mapped against the five hypotheses \emph{before} any
model was fit to real data: the map (\texttt{docs/PRIOR-ART.md}, SHA-256
\texttt{3f36fb62\ldots}) was committed on 2026-08-26 and hash-stamped, and
the hypothesis-status clarifications it forced were registered the same day
as Addendum~3 (Section~\ref{sec:design-prereg}). The verdicts below are
therefore dated evidence of what was known when the design was frozen, not
a reading of the results.

\subsection{Stochastic mortality models and the Cairns--Blake--Dowd programme}

\citet{lee1992modeling} model the log central death rate as
$\log m_{x,t} = \alpha_x + \beta_x \kappa_t$, estimate the bilinear term by
singular value decomposition and extrapolate $\kappa_t$ by a random walk with
drift; the prediction intervals for $\kappa_t$, and hence for life
expectancy, come with the model. \citet{girosi2007understanding} show that
the method is a multivariate random walk with drift on the vector of log
rates with a fixed age profile, so that all of its forecast uncertainty is
$\kappa_t$ uncertainty scaled by $\beta_x$, and that forecast age profiles
become non-smooth; both points matter for reading a Lee--Carter interval.
\citet{brouhns2002poisson} replace the least-squares fit by a Poisson
likelihood with the exposure as offset, the error law appropriate to death
counts, and \citet{brouhns2005bootstrapping} add a semiparametric bootstrap
that carries parameter uncertainty into the projection. Coherent
multi-population variants follow \citet{li2005coherent}, and
\citet{li2017coherent} replace the factor structure by a sparse vector
autoregression on mortality improvements.

\citet{cairns2006two} introduce the two-factor CBD model for post-60
mortality with explicit parameter uncertainty; \citet{renshaw2006cohort} add
a cohort term to Lee--Carter. \citet{cairns2009quantitative} compare eight
such models (M1--M8) on England and Wales and United States data by
Bayesian information criterion, robustness to the fitting period and
biological reasonableness; \citet{cairns2011mortality} take six of them
forward as \emph{density} forecasters and judge their fan charts ex ante;
\citet{dowd2010evaluating} test the residuals of the same six for goodness
of fit; and \citet{dowd2010backtesting} set out a backtesting framework for
their multi-period-ahead density forecasts and apply it ex post to English
and Welsh males. That last paper is the direct ancestor of the present
audit and deserves a precise description. It backtests six models---M1,
M2B, M3B, M5, M6 and M7, each in a parameter-certain and a
parameter-uncertain variant---by locating realised mortality within the
prediction intervals of forecasts made from earlier origins, one population
and one age range at a time. Three of its findings frame ours. It reports
that ``forecast performance tends to improve with higher ages'', the
ex-ante alternative to \Hfour. It notes, in a footnote, that ``the
positions of the individual observations within the prediction intervals
are not independent'' across horizons, and then measures marginal
positions anyway; the joint-path coverage of \Hthree\ is the quantity that
footnote points at and does not compute. And it lists life expectancy,
survival rates and annuity prices as metrics on which ``in principle,
backtests could be conducted'' and confines itself to the mortality rate;
\Hfive\ is that extension. Later members of the programme include the
Bayesian two-population model of \citet{cairns2011bayesian}, the general
procedure for model construction of \citet{hunt2014general} and the CBDX
workhorse of \citet{dowd2020cbdx}. \citet{booth2008mortality} review the
wider field. The \pkg{StMoMo} package \citep{villegas2018stmomo} implements
the family in a common generalised-linear framework and is the numerical
oracle for the classical arms of this study.

Two features of this programme frame ours. It established that stochastic
mortality models are to be judged as density forecasters, not point
forecasters. And its evaluation, apart from \citet{dowd2010backtesting}, was
in-sample or ex ante, on one or two populations, before neural models
existed and before any break in the modern data.

\subsection{Neural extensions of Lee--Carter}

\citet{hainaut2018neural} is an early neural analyser of mortality surfaces.
\citet{richman2021neural} embed country and sex into a network that
generalises the Lee--Carter structure to all HMD populations at once (train
1950--1999, test 2000--2016) and evaluate on mean squared error.
\citet{nigri2019deep} replace the random walk on $\kappa_t$ by an LSTM;
\citet{petnehazi2019mortality} ask whether recurrent networks beat
Lee--Carter, by point error. \citet{perla2021time} show that a shallow
convolutional network can be read as a generalisation of Lee--Carter and
that it transfers from the HMD to the United States Mortality Database.
\citet{scognamiglio2022calibrating} fit Lee--Carter and Poisson Lee--Carter
jointly across populations with a network that replicates their structure.
\citet{marino2023neural} deepen Lee--Carter and add interval estimation
for the future $\kappa_t$ by a bootstrap ensemble of LSTMs;
\citet{miyata2022extending} replace the two-stage procedure by a
variational autoencoder fit in one stage and obtain intervals without MCMC.

\citet{schnurch2022point} is the one neural mortality paper whose title
promises interval forecasts, and it is the closest prior work to ours. On
54 HMD populations at ages 60--89, trained to 2006 and tested on
2007--2016, they report the prediction-interval coverage probability (PICP)
and mean width at a nominal 95\% for Lee--Carter fit on ten and on twenty
years, an autoregressive-cohort model, and recurrent, feed-forward and
convolutional networks: coverages of 74.0, 74.3, 77.2, 86.0, 97.0 and 98.0
per cent respectively, with the neural intervals the widest. Their
intervals come from the variation across independently seeded trainings of
the same network, and their pass rule is one-sided (``minimize MPIW under
the constraint that PICP is at or above'' the nominal level), so
over-coverage is a pass. Two features of that study define what ours adds.
First, the uncertainty mechanism is held fixed by design: the authors
``rule out several existing methods from the literature because they would
require substantial structural changes, such as the insertion of dropout
layers for Monte Carlo dropout'', because doing so ``might decrease the
forecasting performance, the optimization of which is still our main
goal''. The 74-versus-98 per cent gap is therefore confounded between
architecture and mechanism, which is exactly what the crossed design of
Section~\ref{sec:design} separates. Second, their classical arms see ten or
twenty training years while the networks see every year of all 54
populations, so part of the classical under-coverage is a training-window
artefact; the expanding-origin rule of Section~\ref{sec:design-regimes}
gives every family the same window. Their Table~2 also shows the
point-versus-proper-score inversion of \Hone\ (mean squared error ranks the
networks first; Poisson deviance ranks Lee--Carter first), which they
attribute to exposure weighting rather than to distributional sharpness.

Adjacent non-neural machine-learning comparators are the gradient-boosted
multi-population model of \citet{miao2025gradient}, the tree-based analysis
of \citet{deprez2017machine}, the pension-fund application of
\citet{demelo2025actuarial}, and the Markov-switching Bayesian VAR on an
age-partitioned Lee--Carter of \citet{fu2025applying}, which is a Bayesian
rather than a neural route to uncertainty under structural change.

The evaluation profile of this lineage is documented in
Section~\ref{sec:intro}: point accuracy throughout; intervals plotted where
they exist; coverage computed without a sharpness penalty where it is
computed at all; proper scores rarely. Its stable-regime coverage record is
nevertheless informative, and it already settles part of \Htwo: in periods
without a break, classical intervals under-cover and bootstrap-ensemble
neural intervals over-cover \citep{schnurch2022point,marino2023neural}. The
registered ``under-cover for every family'' clause was therefore falsified
in the stable regime before our first fit, and the stable regime of this
paper is a replication arm (Section~\ref{sec:design-status}). The four
neural families in our grid (Section~\ref{sec:design-families}) are drawn
from this lineage.

\subsection{Bayesian and Gaussian-process mortality models}

\citet{ludkovski2018gaussian} model mortality surfaces by Gaussian-process
regression; \citet{huynh2021multi} extend this to multi-output GPs over many
HMD populations and, alone in our library, evaluate by CRPS as the score is
meant to be used. \citet{barigou2023bayesian} average Bayesian mortality
models by leave-future-out validation with the expected log predictive
density and CRPS across four countries, and find that ``the `optimal' model
depends on the country and the scoring rule studied'';
\citet{goes2025bayesian} report, on an England and Wales war-shock window,
a top-one inversion between CRPS and the point and log-score rankings.
These two papers, with \citet{schnurch2022point}, are why \Hone\ is prior
art and is carried as a harness-validity check rather than a finding
(Section~\ref{sec:design-status}). \citet{liu2020bayesian} fit a Bayesian
Poisson log-normal model with regularised time structure. These are the
models whose predictive distributions are honest by construction, and their
evaluation practice is the one we adopt for every family; a multi-output GP
is one of the ten families audited.

\subsection{COVID-era mortality modelling}

The pandemic literature in actuarial science is about \emph{estimation with}
the pandemic years, not \emph{evaluation across} them.
\citet{cairns2020impact} frame COVID-19 deaths as accelerated deaths with
implications for higher-age mortality; \citet{robben2022assessing} quantify
what one pandemic data point does to the calibration and projections of a
stochastic multi-population model; \citet{goes2025bayesian} extend
Lee--Carter with a vanishing jump so that pandemic years neither dominate the
trend nor have to be deleted; \citet{robben2024catastrophe} add catastrophe
risk to a multi-population model; \citet{vanberkum2025estimating} extend the
Li--Lee model to measure the impact on granular data;
\citet{nalmpatian2023advanced} combine generalised additive models with
machine learning to estimate the trend deviation; and
\citet{robben2025granular} model epidemic and temperature shocks by regime
switching. \citet{barigou2023bayesian}, as noted, simulate a pandemic
perturbation because real post-shock data did not yet exist, and remark
that adapting their validation to ``a change of regime'' is ``beyond the
scope'' of their paper. \citet{goes2025bayesian} state the constraint
directly: because the pandemic ``predominantly impacts the most recent
years of data'', an out-of-sample test was ``impractical'' and their COVID
analysis ``focuses solely on evaluating the in-sample fit''. None of these
papers evaluates prediction intervals that were fit before 2020 against
what happened in 2020--2024; with twenty populations now final through
2024, the constraint that stopped them is gone. That is the gap this audit
fills, and the papers above are what a modeller would reach for afterwards
to repair a failure we document; repairing it is not our contribution.

\subsection{Proper scoring rules, calibration and sharpness}

A scoring rule is proper if the expected score is optimised by the true
distribution and strictly proper if uniquely so; the logarithmic score and
the continuous ranked probability score are the two canonical examples
\citep{gneiting2007strictly}, and the interval (Winkler) score is the proper
score for a central prediction interval \citep{winkler1972decision}.
\citet{gneiting2007probabilistic} set out the paradigm that the present audit
adopts wholesale---maximise sharpness subject to calibration---together
with the probability integral transform (PIT) as the diagnostic of
probabilistic calibration, following the prequential principle of
\citet{dawid1984present}. \citet{czado2009predictive} give the randomised
PIT for count data, which we need because deaths are counts and predictive
samples tie. The decomposition of a proper score into uncertainty,
resolution and miscalibration goes back to \citet{murphy1973new};
\citet{pohle2020murphy} generalises it to all forecast types and
\citet{dimitriadis2021stable} give a stable estimator via isotonic
recalibration. \citet{dawid2014theory} survey the theory;
\citet{yanagisawa2023proper} extend proper scoring to survival-time
distributions, which is the natural target when the object of interest is a
lifetime rather than a rate.

\subsection{Conformal prediction}

Conformal prediction \citep{vovk2005algorithmic} produces prediction sets
with finite-sample marginal coverage under exchangeability of the calibration
and test residuals; split conformal \citep{lei2018distribution} makes it
cheap for any black-box regressor, and conformalised quantile regression
\citep{romano2019conformalized} makes the sets adaptive. Its relevance to a
time series is the question of what happens when exchangeability fails.
\citet{tibshirani2019conformal} reweight under covariate shift;
\citet{barber2023conformal} bound the coverage gap of weighted conformal
sets by a total-variation distance between calibration and test
distributions---a bound that a structural break is designed to make large.
\citet{xu2021conformal} propose ensemble batch prediction intervals (EnbPI)
for dynamic time series without exchangeability; \citet{gibbs2021adaptive}
adapt the miscoverage level online under distribution shift, and
\citet{gibbs2024conformal} extend this to arbitrary shifts. For multi-step
forecasts, where the object of interest is the whole path, joint coverage
is solved prior art in machine learning: \citet{stankeviciute2021conformal}
obtain simultaneous multi-horizon bands from a recurrent network by a
Bonferroni allocation, \citet{sun2024copula} calibrate a joint band through
a copula of the per-horizon scores, \citet{messoudi2021copula} build
copula-based multi-target sets, and \citet{diquigiovanni2022conformal} give
functional bands from a scaled supremum-norm score. \citet{sun2025conformal}
address change points directly. \citet{angelopoulos2023gentle} is the
standard primer.

We found no application of conformal prediction to mortality forecasting:
in the 52-paper corpus of the prior-art map, ``conformal'' has no hit in any
mortality text, and neither do ``Monte Carlo dropout'' or ``deep
ensemble''. Three conformal mechanisms enter our grid as \emph{audited
arms}, on the same footing as the native and ensemble mechanisms; the
joint-path arm is positioned against \citet{stankeviciute2021conformal}
and \citet{diquigiovanni2022conformal}, whose construction it uses, and we
claim the first \emph{measurement} of the marginal-to-joint coverage gap in
mortality forecasting, not the method. We make no claim that any conformal
arm fixes anything, and we deliberately do not invoke the finite-sample
guarantee under a break, for exactly the reason \citet{barber2023conformal}
give.

\subsection{Uncertainty mechanisms for neural networks}

Deep ensembles \citep{lakshminarayanan2017simple} estimate predictive
uncertainty from the disagreement of independently initialised networks;
Monte Carlo dropout \citep{gal2016dropout} reads dropout at prediction time
as approximate Bayesian inference. Both capture variability that the
training data can reveal; neither is designed to anticipate a regime the
training data never contained. Whether that matters for coverage at a break
is one of the questions the crossed design answers.

\subsection{Forecast-comparison inference}

\citet{diebold1995comparing} test equal predictive accuracy on a
loss-differential series; \citet{hansen2011model} construct the model
confidence set, the set of models not significantly outperformed at a given
confidence level; \citet{shang2018model} apply the latter to age-specific
mortality (Japanese national and prefectural data, retirement ages), the one
prior use we found in this setting. Because our shift regime is a single
common shock observed in twenty correlated populations, the effective
number of independent observations is nearer one than twenty, and ordinary
Diebold--Mariano standard errors overstate the information available; we
cluster by population, impose the null, and use the wild cluster bootstrap
of \citet{cameron2008bootstrap} with the six-point weights of
\citet{webb2023reworking}, which \citet{mackinnon2017wild} show remain
reliable at the small number of clusters we have. Rolling-origin evaluation
follows \citet{tashman2000out}; pre-registration follows the practice
advocated by \citet{nosek2018preregistration}.

\subsection{What is and is not prior art}

For the record, and because it was written down before any result existed,
the prior-art map's verdict on each hypothesis is: \Hone\ is prior art
\citep{barigou2023bayesian,goes2025bayesian,schnurch2022point}; the
stable-regime half of \Htwo\ is prior art and contradicts the registered
universality, while its shift-regime half is untested anywhere; \Hthree\ is
novel in mortality as a measurement and not novel as a method; \Hfour's
registered direction is contradicted by \citet{dowd2010backtesting}; and
\Hfive's construction is prior art while an ex-post coverage audit of the
derived quantities is not found. The white space that remains, and that
this paper occupies, is the crossed family-by-mechanism design, the
conformal arms, the marginal-to-joint gap, the derived-quantity coverage
audit, and the out-of-sample pandemic test the field declined for lack of
data. Two items are still unread at the time of writing and are named so
that their absence is visible: the online supplement of
\citet{schnurch2022point} (Section~C.3, implied annuity present values),
which is the one live threat to \Hfive's residual novelty if it audits
annuity-interval coverage rather than presenting point values, and the full
text of \citet{stankeviciute2021conformal}, cited above from its
description in later work.

\section{Data}
\label{sec:data}

\subsection{Source and vintage}

All data are from the Human Mortality Database \citep{hmd2026}, bulk
``zipped data files'' download, files \texttt{Deaths\_1x1},
\texttt{Exposures\_1x1}, \texttt{Mx\_1x1} and the period life tables, last
modified 15~June~2026, release DOI 10.4054/HMD.Countries.20260615, and
constructed under Methods Protocol version~6, revision of 5~August~2025
\citep{wilmoth2025methods}. HMD has issued a citable DOI per release since
September 2025; the DOI, the download date (2026-08-25) and the SHA-256
checksums of every input file in \texttt{data/MANIFEST.sha256} together pin
the exact bytes. The data themselves are not redistributed, in keeping with
the HMD user agreement, and any table in this paper can be reproduced from
a fresh download that matches the manifest. The bulk deaths and exposures
files each carry fifty distinct population codes (HMD's \texttt{PopName}),
the count recorded in the manifest header: forty-one countries or areas,
plus nine further codes for the overlapping sub-population and civilian
variants that HMD publishes for Germany (two), France (one), the United
Kingdom (four) and New Zealand (two). This is the first HMD vintage in which
twenty of them carry \emph{final} single-age annual data through 2024,
which is what makes the shift regime (Section~\ref{sec:design}) possible at
$h=5$; one (DNK) already reaches 2025, which the shift regime does not use.

Only the deaths and exposures enter the models. The \texttt{Mx\_1x1} file is
used once, to validate the parser: our $D/E$ reproduces the published $m_x$
exactly at spot checks. The period life tables are used once, for validation
gate~3 (Section~\ref{sec:design-gates}).

\subsection{Cells, ages and exposure}

A cell is a (population, sex, calendar year, single age) Lexis square with
death count $D_{x,t}$ and central exposure $E_{x,t}$. HMD reports ages
$0,\dots,109$ and an open group $110{+}$. We model ages $0$--$99$ and drop
the rest, for three reasons. First, the published life tables smooth death
rates at the oldest ages (a Kannisto fit above age~80 under Protocol~v6), so
the old-age surface is partly model output rather than observation. Second,
the open group needs a closure assumption that no forecasting family in the
grid shares. Third, exposure vanishes: read with the study's own parser
(\texttt{mortcal.data.hmd.read\_bulk\_1x1}), the vintage contains 24{,}083
population--year--age rows with $E_{x,t}=0$ for at least one sex (15{,}494
for both sexes), almost all at age~90 or above and concentrated in the
small populations; a Poisson model with offset $\log E_{x,t}$ is undefined
there. Sexes are modelled separately throughout; the \emph{Total} column is
never used.

Inside $0$--$99$ a zero-exposure cell is a \emph{structural} zero, not a
missing value: in every such cell in the shift training panel HMD's
1~January population count is zero at both corners of the Lexis square, so
nobody was alive at that age and $D_{x,t}=0$ as well. Under the Poisson
likelihood with log-exposure offset such a cell has mean $E\exp(\eta)=0$
and observation $0$: its log-likelihood contribution, score and Fisher
information are all identically zero, so the maximum-likelihood estimate
with the cell included equals the estimate under indicator weights
\emph{exactly}; only the numerical evaluation of $0\log 0$ fails. Every
family therefore carries the cell-weight matrix $W_{x,t}=\ind\{E_{x,t}>0\}$
(the \texttt{wxt} convention of \pkg{StMoMo}), which Addendum~3~\S1 of the
pre-registration records as the numerically correct evaluation of the
registered objective rather than as a deviation; the weighted code paths
reproduce the unweighted fits bit-identically on panels without zeros,
which is tested. The shift training panel contains 99 such cells, all at
ages 95--99, in seven population--sex series (CHE, FIN, ISL, LUX and SWE
males, CHE and ISL females), every one before 1970. In a \emph{test}
window a zero-exposure cell has no observation to score; the scorer masks
an age only if every horizon's observation is finite, records the number of
ages scored, and computes life-table functionals on the maximal contiguous
scored age range from age~0 on both the predictive and the observed side
(Addendum~3~\S3). This binds only in the placebo regime, at ages 98--99, in
a handful of Swiss, Finnish and Icelandic cells. Zero-\emph{death} cells
with positive exposure are ordinary observations under the Poisson
likelihood and remain in every fit.

\subsection{Fractional deaths, half-count rates and the rounding convention}
\label{sec:data-rounding}

HMD death counts are not integers. Deaths are allocated to Lexis triangles
and then re-aggregated, and deaths of unknown age are redistributed, so a
cell may record $3842.31$ deaths; this is the norm, not an anomaly. It has
two consequences that the pre-registration and its addenda fix in advance.

First, a Poisson log score evaluated at a non-integer count is undefined.
We score the Poisson predictive at the death count rounded to the nearest
integer with halves rounded \emph{up}, $\tilde D_{x,t} = \lfloor D_{x,t} +
\tfrac12 \rfloor$, and report as a registered sensitivity the CRPS of the
predictive count distribution against the unrounded $D_{x,t}$, which needs
no rounding. The half-up rule is implemented as written; IEEE
round-half-to-even was rejected because a small share of HMD cells land
exactly on a half and its direction would then depend on the parity of the
integer part, a cell-dependent perturbation no reader would infer from
``nearest integer''. Nine of the twenty shift populations carry
non-integer counts, two of them (JPN, USA) in essentially every cell; the
other eleven are exactly integer, and the per-population non-integer share
is reported so that the rounding's reach is visible.

Second, a cell with zero observed deaths has no finite log rate. Rate-scale
quantities---observed and predictive alike---use the half-count continuity
correction of demographic practice,
\begin{equation}
  y_{x,t} = \log\frac{\max(\tilde D_{x,t},\,\tfrac12)}{E_{x,t}},
  \label{eq:halfcount}
\end{equation}
so that both sides of every rate-scale score live on the same lattice
(Addendum~2~\S2, Addendum~3~\S5). Cells are never dropped for having no
deaths; the number of test cells in which the correction binds
($D_{x,t}<\tfrac12$, counted on the full age range before any model's age
mask) is reported per population, sex and regime. Such cells are
concentrated in the smallest populations, ISL and LUX, at the youngest and
oldest ages, where they are a material fraction of the 2020--2024 test
window; they are rare elsewhere. The same convention is applied to the
conformal calibration residuals (Section~\ref{sec:design-mechanisms}), where
a single zero-death calibration cell under a small rate floor would
otherwise inflate a band radius by an order of magnitude.

One deliberate exception: the \emph{observed} period life table used for
the realised $e_0$, $e_{65}$ and $\annuity$ of \Hfive\ keeps the raw
fractional deaths and a negligible rate floor, because at an age where
nobody died the crude rate $m_x=0$ is the correct life-table statement,
while a half death would shift the observed $e_0$ of a small population by
more than validation gate~3 tolerates. The two conventions---half a death
on the rate scale, no death in the observed life table, for the same
cell---are both registered (Addendum~3~\S5).

One historical detail of the triangle split matters for the placebo regime.
HMD's linear model for apportioning $1\times1$ death squares into Lexis
triangles (Protocol~v6, Appendix~A) carries a 1918--1919 influenza correction
extrapolated from Sweden and France. It changes how deaths are shared between
the two triangles of a square, not the $1\times1$ counts we use, and leaves
only a small model-based imprint on the 1918--1919 exposures.

\subsection{The likelihood}

Deaths are modelled as counts: $D_{x,t}\sim\mathrm{Poisson}(E_{x,t}\,m_{x,t})$
with $\log E_{x,t}$ as offset, following \citet{brouhns2002poisson}.
Nothing in the study models rates as Gaussian on the raw scale. The
classical families are fit by this likelihood, with two deliberate
exceptions retained as historical baselines: the original singular-value
Lee--Carter with its least-squares fit on log rates, and the CBD model with
its per-year least-squares fit on the logit of $q_{x,t}$. The neural
families fit the same Poisson deviance with the log-exposure offset
(neural-LC and the shallow CNN), a negative-binomial log-likelihood
(the distributional head), or, for the LSTM on $\kappa_t$, inherit the
SVD Lee--Carter stage and fit the index by squared error, as their
reference papers do; the multi-output Gaussian process places a Gaussian
likelihood on the half-count log rate~\eqref{eq:halfcount} with a
Poisson-level noise floor. The loss of every family is written in the
family specification committed before the neural sweep
(\texttt{docs/NEURAL-SPEC.md}) and summarised in
Section~\ref{sec:design-families}.

\subsection{Populations by regime}

Population codes are HMD's. The sets were fixed in the pre-registration from
an audit of the files on disk, not from results.

\paragraph{Shift (primary), 20 populations.} Those whose final annual data
reach 2024 in this vintage: BEL, CHE, CHL, DNK, EST, FIN, HKG, HRV, ISL, JPN,
KOR, LTU, LUX, LVA, NOR, PRT, SVK, SWE, TWN, USA. Every population trains on
the maximal contiguous block of complete years ending in 2019, from its own
first year (Addendum~3~\S2). The rule binds once: Belgium's 1914--1918
deaths and exposures are missing at every age---1{,}000 modelled cells,
$100$ ages $\times$ $2$ sexes $\times$ $5$ years, the only genuinely missing
data in the shift training panel (Addendum~3~\S2)---so BEL trains on
1919--2019; without the
rule a random-walk drift would treat 1913 to 1919 as one step and a
positional cohort index would mislabel every cohort after 1913 by five
years. The series lengths at the 2019 origin are very unequal, from
seventeen years (KOR) and nineteen (HRV) to more than two centuries (SWE);
this matters for which uncertainty mechanisms can be constructed on the
shortest panels (Section~\ref{sec:design-grid}).

\paragraph{Stable (control).} The same twenty populations, each from its own
first year of data, at expanding origins $T=1990,1992,\dots,2014$ with test
years $T+1,\dots,T+5$ capped at 2019. An (origin, population, sex) triple is
admissible only with at least fifteen complete training years
(Addendum~3~\S4): shorter windows produce plausible-looking junk, such as a
random-walk innovation variance estimated from two differences. The stable
panel is therefore unbalanced---the number of buildable population--sex
series grows across origins as the late-starting Baltic, East Asian and
South American series pass the floor---and the effective cluster count per
origin is reported alongside every stable-regime summary.

\paragraph{Placebo break, 11 populations.} Those with continuous coverage
over 1908--1922 and a start year no later than 1900: CHE, DNK, FIN, FRATNP,
GBRTENW, GBR\_SCO, ISL, ITA, NLD, NOR, SWE. Belgium is excluded because its
occupation-era data for 1914--1918 are missing at every age: 1{,}000
modelled cells in the registered unit of Addendum~3~\S2 ($100$ ages
$\times$ $2$ sexes $\times$ $5$ years), which appear in the raw $1\times1$
file as 555 rows (111 ages including $110{+}$, one row per age-year, 5
years) that are present but carry HMD's ``.'' missing marker in the Female,
Male and Total columns; HMD's own documentation warns of the gap. FRACNP and GBRCENW are excluded as
civilian-population variants overlapping FRATNP and GBRTENW; keeping both
would double-count one country, and HMD additionally flags FRACNP's 1919
denominators as implausible, so it is not used even as a sensitivity.

Addendum~1 to the pre-registration (2026-08-26, still before any real-data
fit) stratifies the eleven from HMD's country documentation into
\emph{neutral, register-based} series whose only break in the window is the
1918 influenza pandemic (CHE, DNK, FIN, ISL, NLD, NOR, SWE);
\emph{belligerent total-population} series in which HMD reconstructs
military deaths at male ages of roughly 18--45 (FRATNP, GBRTENW, ITA); and one
\emph{civilian-only} series in which deaths abroad are excluded by
construction (GBR\_SCO). The neutral stratum is the clean analogue for the
1918-versus-2020 comparison; the pooled placebo result is always reported
alongside the strata.

\begin{table}[htbp]\centering
\caption{Populations by regime: first and last year of single-age data in
the pinned vintage, training years available at the shift origin,
zero-exposure cells at ages $0$--$99$, and zero-death test cells
($D<\tfrac12$) per sex in the shift window. Generated from the data files,
not from any model output.}
\label{tab:populations}
\inputtable{tab-populations}
\end{table}

\subsection{Data-quality caveats carried into the analysis}

\begin{itemize}
\item \emph{Exposure conventions.} Exposures are estimated from census or
  register populations and births, with conventions that differ by country
  and era; \citet{cairns2016phantoms} show how cohort-size errors of several
  percent arise from uneven historical births. We take HMD exposures as
  given and note that any such error is common to every model in a cell.
\item \emph{The placebo window.} Pre-1950 data are patchier and the war years
  are, for some countries, reconstructed. The public HMD background and
  documentation file of each placebo population, the Methods Protocol and
  the release notes were read before any placebo result
  (\texttt{docs/DATA-PREREQS.md}); the per-country ``Notes'' files behind
  the HMD login were not consulted. No population carries an explicit
  1914--1922 quality flag, and the concerns that do exist are the strata
  above plus three specific items that become pre-specified sensitivities:
  GBR\_SCO's 1912--1938 population is interpolated by a spline HMD itself
  calls unreliable (sensitivity: drop GBR\_SCO); Denmark gained South Jutland
  in 1921, inside the test window (sensitivity: DNK with and without
  1921--1922); and the neutral stratum alone is reported as the cleanest
  pandemic-only comparison.
\item \emph{Revisability of 2024.} HMD carries no per-year provisional flag;
  each release re-issues whole series, so ``final through 2024'' means final
  as of this release. Known denominator risks in the shift panel: USA
  2020--2025 exposures blend preliminary 2020-census inputs; CHL exposures
  rest on the 2017 census pending revision; TWN 2024 deaths are provisional
  according to the short-term mortality fluctuations series. Register-based
  populations (the Nordic countries, CHE, EST, LVA, LTU) are effectively
  final. The pre-registered robustness run drops 2024 and re-scores the shift
  regime at $h\le4$; Addendum~1 adds dropping USA and CHL, a
  register-versus-census contrast, and a re-run on the next HMD vintage if
  one is released before submission.
\item \emph{Differing observation windows.} Series start between the
  mid-eighteenth century and the 2000s; the stable regime is therefore an
  unbalanced panel over origins, and the shortest shift panels cannot
  support every mechanism. This is why every inference is clustered by
  population rather than pooled over cells, and why every contrast is
  computed on the cells common to the arms it compares
  (Section~\ref{sec:design-grid}).
\end{itemize}

\subsection{Life-table conventions}

Life expectancy and annuity factors are computed from every predictive
sample's $m_x$ vector by one period-life-table routine:
$q_x = m_x/(1+(1-a_x)m_x)$; $a_0 = 0.07+1.7\,m_0$ capped to $[0.01,0.35]$,
a linear infant-separation rule in the spirit of \citet{andreev2015average};
$a_x=\tfrac12$ elsewhere; and constant-hazard closure at the last age,
$e_A = 1/m_A$. HMD itself uses the piecewise Andreev--Kingkade coefficients
and its own old-age treatment, so the two tables are not identical;
validation gate~3 requires our $e_0$ and $e_{65}$ on HMD's published $m_x$
to reproduce HMD's published values within 0.15 years, which is the size of
the documented $a_x$ simplifications. Annuity factors follow the standard
life-contingency definition \citep{dickson2020actuarial} on the period table
treated as static (Section~\ref{sec:metrics-actuarial}). A family defined
only on part of the age range (CBD, ages 55--99) yields a table truncated at
its first age: $e_0$ is then undefined and $e_{65}$ and $\annuity$ are
computed on the truncated table, on the predictive and observed sides
alike, with the age range recorded.

\section{Study design}
\label{sec:design}

\subsection{Pre-registration}

The design was written down, committed and hash-stamped on 2026-08-25,
before any model was fit to real HMD data; its SHA-256 digest is
\begin{center}
\small\texttt{bbbe3860a5446d063e186e56555afa893b07c36cf15a990d07567471affe50be}
\end{center}
The registration fixes the analytic choices: regimes and their exact years,
population lists, the model families and mechanisms, the admissible grid,
every metric and its convention, the inference procedure, the validation
gates, and the five hypotheses. It does not, and could not, hide the outcome
of 2020--2024 from us; pre-registration here protects against analytic
flexibility, not against hypothesis surprise. One ordering deviation is
recorded honestly: the registration was frozen before
\citet{dowd2010backtesting} had been read. Nothing in the frozen design
depends on that paper's content; what it can change is the framing of
novelty in the text, which is not frozen. One addendum has been registered
(Addendum~1, 2026-08-26, still before any real-data fit): it adds the
placebo strata and the denominator sensitivities described in
Section~\ref{sec:data}, and changes nothing in the hypotheses, populations,
splits, horizons, metrics or inference plan. \todo{cite the commit hash of
Addendum~1} Every subsequent deviation is reported as a deviation in
Section~\ref{sec:design-deviations}.

\subsection{Three regimes}
\label{sec:design-regimes}

\begin{table}[htbp]\centering
\caption{The three pre-registered evaluation regimes.}
\label{tab:regimes}
\small
\begin{tabular}{@{}lllcc@{}}
\toprule
Regime & Training years & Test years & $h$ & Populations \\
\midrule
Stable (control) & $\le T$, $T=1990,1992,\dots,2014$ & $T{+}1,\dots,T{+}5$ ($\le 2019$) & $1$--$5$ & 20 \\
Shift (primary)  & $\le 2019$ & 2020--2024 & $1$--$5$ & 20 \\
Placebo break    & $\le 1913$ & 1914--1922 & $1$--$9$ & 11 \\
\bottomrule
\end{tabular}
\end{table}

The \emph{stable} regime measures calibration where nothing breaks; it is
the reference level against which \Htwo\ and \Hfive\ are stated. The
\emph{shift} regime is the audit. The \emph{placebo} regime is a second
event---the First World War and the 1918 influenza pandemic---with a longer
horizon, used descriptively: a side-by-side of 1914--1922 and 2020--2024
coverage answers ``is this just COVID?'' without a transfer regression.
Every window is expanding-origin; training always uses all years up to the
origin from the population's own first year. A random split is never used
anywhere, because mortality panels leak trivially under random splitting:
adjacent ages and years are near-duplicates.

\subsection{Model families}
\label{sec:design-families}

Ten families, chosen to span the classical, statistical-learning and neural
lineages of Section~\ref{sec:related}. Each is specified by its reference
paper; the repository documents any deviation.

\begin{enumerate}
\item \textbf{Lee--Carter (SVD).} $\log m_{x,t}=\alpha_x+\beta_x\kappa_t$,
  SVD fit with $\sum_x\beta_x=1$, $\sum_t\kappa_t=0$; random walk with drift
  on $\kappa_t$ including drift-estimation uncertainty
  \citep{lee1992modeling}.
\item \textbf{Poisson Lee--Carter.} Same structure, Poisson likelihood with
  log-exposure offset, fit by Newton's method \citep{brouhns2002poisson}.
\item \textbf{CBD (M5).} $\operatorname{logit} q_{x,t} = \kappa^{(1)}_t +
  \kappa^{(2)}_t (x-\bar x)$ on ages 55--99, bivariate random walk with drift
  \citep{cairns2006two}.
\item \textbf{APC / Renshaw--Haberman (M2-A).} Lee--Carter with a cohort term
  $\gamma_{t-x}$ under the identifiability constraints of
  \citet{cairns2009quantitative}; ARIMA projection of the cohort index.
  \citep{renshaw2006cohort}.
\item \textbf{Sparse VAR.} Banded vector autoregression on log-rate
  improvements with a penalised fit \citep{li2017coherent}.
\item \textbf{Multi-output Gaussian process.} GP regression over (age, year)
  with a shared-across-populations kernel structure
  \citep{huynh2021multi}. \todo{final kernel and inducing-point choices;
  document exclusion if the GPU budget forbids}
\item \textbf{Neural-LC.} Country and sex embeddings feeding a network that
  generalises the Lee--Carter surface \citep{richman2021neural}.
\item \textbf{Shallow CNN-LC.} Convolutional generalisation of Lee--Carter
  \citep{perla2021time}.
\item \textbf{LSTM-on-$\kappa_t$.} Lee--Carter age structure with a recurrent
  network in place of the random walk \citep{nigri2019deep}.
\item \textbf{Distributional Poisson/NB head.} A network whose output layer
  parameterises a Poisson or negative-binomial count distribution for
  $D_{x,t}$ given $\log E_{x,t}$; the only neural family with a model-native
  predictive distribution.
\end{enumerate}

Hyperparameters of the neural families are tuned on an inner \emph{time}
split---the final five pre-cutoff years of each training window---never on a
test year and never on a random subset.

\subsection{Uncertainty mechanisms}
\label{sec:design-mechanisms}

Seven mechanisms, each a way of turning a fitted family into predictive
samples.

\begin{enumerate}
\item \textbf{Model-native.} The predictive distribution the reference paper
  itself supplies: random-walk innovation and drift uncertainty for
  Lee--Carter-type models, the GP posterior, the count distribution of the
  distributional head.
\item \textbf{Semiparametric Poisson bootstrap.} Refit the family $B$ times to
  pseudo-deaths $D^{(b)}\sim\mathrm{Poisson}(E\hat m)$ and pool the refits'
  native samples \citep{brouhns2005bootstrapping}; adds parameter uncertainty
  the native mechanism conditions away.
\item \textbf{Deep ensemble.} $M=10$ independently seeded trainings, samples
  pooled with equal weight \citep{lakshminarayanan2017simple}.
\item \textbf{Monte Carlo dropout.} Dropout active at prediction time,
  stochastic forward passes as samples \citep{gal2016dropout}.
\item \textbf{Split conformal.} Absolute-residual conformal intervals on the
  log-rate scale \citep{lei2018distribution}, Mondrian-stratified by age
  band ($0$--$24$, $25$--$64$, $65$--$99$) and by horizon
  \citep{vovk2005algorithmic}.
\item \textbf{EnbPI.} Ensemble batch prediction intervals adapted from the
  online original to annual multi-step forecasting \citep{xu2021conformal};
  the adaptation is documented in the repository.
\item \textbf{Copula / joint-path conformal.} A simultaneous band over
  horizons from a scaled sup-norm nonconformity score
  \citep{diquigiovanni2022conformal}, in the spirit of
  \citet{sun2024copula} and \citet{messoudi2021copula}; the only mechanism
  whose band is calibrated to be joint over $h=1,\dots,H$.
\end{enumerate}

Two disciplines apply to every conformal arm. The calibration residuals are
drawn only from years at or before the training cutoff, on an inner time
split in which the calibration years are later than the years the base fit
is trained on; no mechanism sees a test year. And a conformal wrapper yields
an interval, not a distribution: to preserve the single sample interface the
harness draws uniformly inside each cell's interval, independently across
cells. That convention keeps coverage, width and interval score exact at the
construction level, but it makes CRPS, log score and PIT for conformal cells
a placeholder rather than a distributional claim. Those scores are computed,
flagged, and reported only in Appendix~\ref{app:secondary}; they are never
ranked against distributional mechanisms.

\subsection{The crossed grid}
\label{sec:design-grid}

Not every family admits every mechanism. Deterministic classical fits have
no seed variance, so ensembles and dropout are undefined for them; the GP
posterior already integrates parameter uncertainty, so a bootstrap would be
redundant; and neural networks have no closed-form native predictive except
through the distributional head. Table~\ref{tab:grid} enumerates the
admissible cells: 50 primary and a handful of secondary ones. The three
conformal columns are wrappers around the same underlying fits as the native
and ensemble columns, not refits, so the expensive axis is neural families
$\times$ ensemble members $\times$ expanding origins.

\begin{table}[htbp]\centering
\caption{Admissible model-family $\times$ uncertainty-mechanism grid, transcribed
from the pre-registration.}
\label{tab:grid}
\inputtable{tab-grid}
\end{table}

Because the grid is ragged, no claim in this paper is a full-factorial main
effect. Every comparison is a \emph{contrast within an admissible sub-grid},
and the pre-registration names the contrasts that carry the hypotheses:
conformal versus native within the classical families; ensemble versus
dropout within the neural families; bootstrap versus native within the
classical families; and joint-path conformal versus split conformal across
all families.

\subsection{One interface, one scoring path}
\label{sec:design-interface}

Every family--mechanism cell implements the same two calls,
\[
  \texttt{fit}(D, E) \quad\text{and}\quad
  \texttt{sample\_mx}(h, n, \text{rng}) \;\to\; [n, h, n_{\text{ages}}],
\]
returning $n$ pathwise trajectories of $m_{x,T+1},\dots,m_{x,T+h}$. Point
forecasts, intervals, proper scores, PIT values, joint path coverage and the
life-table functionals are all computed from these samples by one evaluation
module. No model has a bespoke metric implementation, so no model can be
accidentally flattered by one. The point forecast for RMSE and MAE is the
pointwise median of the log-rate samples. \todo{unify the point functional
for $e_0$ (currently the sample mean) with the median used for log rates}

\subsection{Seeds}

The global seed is 20260825. Each cell derives its own generator from a seed
sequence over (global seed, origin, population, sex, family, mechanism);
ensemble member $j$ uses global seed plus $j$. Every seed is written to the
results file alongside the metrics.

\subsection{Validation gates}
\label{sec:design-gates}

Three gates had to pass before any real-data result could be read. All three
have passed.

\paragraph{Gate 1: synthetic truth.} Data are simulated from a known Poisson
Lee--Carter process (40 ages, 60 training years, horizons $1$--$10$,
exposures $10^5$, thirty independent worlds, 1{,}500 predictive samples
each). The correctly specified Poisson Lee--Carter must attain nominal
95\% coverage within Monte Carlo tolerance and a near-uniform PIT histogram;
the harness must also detect an injected break as under-coverage. It does
both. If this gate had failed, every coverage number on real data would have
been measuring a bug. One documented finding from the gate: the two-stage
SVD Lee--Carter is mildly \emph{over}-dispersed on Poisson data, because the
SVD $\kappa_t$ carries observation noise into the random-walk innovation
variance; its intervals are slightly wide and its PIT centre-heavy. This is
a property of the method, not of the harness, and it is carried into the
interpretation of the SVD-LC rows.

\paragraph{Gate 2: oracle parity.} The Python Poisson Lee--Carter must match
R \pkg{StMoMo} (version 0.4.1 under R~4.6.1) on a shared subset. On Swedish
males, 1950--2000, ages $0$--$99$, the two fits have identical Poisson
log-likelihood and a maximum relative difference in $\beta_x$ of
$2.4\times10^{-14}$ when the Newton iteration is run to its cap.
\todo{quote the $\alpha_x$ and $\kappa_t$ relative differences from the
parity script} The \pkg{StMoMo} parameter dump is
tracked in the repository so the check can be re-run without R.

\paragraph{Gate 3: life-table parity.} Our period life table on HMD's
published $m_x$ reproduces HMD's published $e_0$ and $e_{65}$ within
0.15~years at spot checks (Section~\ref{sec:data}).

\subsection{Inference plan}

The shift regime is one event. Twenty populations experienced one common
shock, so the effective number of independent observations of ``what a
pandemic does to an interval'' is nearer one than twenty, and nearer one
than the thousands of cells in the panel. Inference is therefore framed as
an event study: population is the cluster, every test resamples clusters
rather than cells, PIT is pooled within cluster before it is compared
across clusters, and the placebo regime is treated as a second event rather
than as more draws from the same one. The procedures are defined in
Section~\ref{sec:metrics-inference}.

\subsection{Deviations from the pre-registration}
\label{sec:design-deviations}

\begin{itemize}
\item Ordering: \citet{dowd2010backtesting} read after the registration was
  frozen (above).
\item Addendum~1 (2026-08-26): pre-specified placebo strata and shift-regime
  denominator sensitivities. Not a deviation---nothing frozen changed---but
  recorded here so the ledger is complete.
\item \todo{append every subsequent deviation here with its date and reason;
  an empty list at submission is a claim, not a default}
\end{itemize}

\section{Metrics}
\label{sec:metrics}

\subsection{Notation}

Fix a population, a sex and an origin $T$. For age $x$ and horizon $h$ the
target on the rate scale is $y_{x,h} = \log m_{x,T+h}$ and on the count scale
it is the observed death count $D_{x,T+h}$ with known exposure $E_{x,T+h}$.
A cell of the grid supplies $n$ predictive sample paths
$z^{(j)}_{x,h}=\log m^{(j)}_{x,T+h}$, $j=1,\dots,n$, jointly over $(x,h)$;
$\hat F_{x,h}$ denotes the empirical predictive distribution of
$\{z^{(j)}_{x,h}\}_j$ and $\hat F^{-1}_{x,h}(p)$ its empirical quantile.
All scores are negatively oriented: smaller is better. Averages over cells
are denoted by an overbar; averaging is always within population--sex
first, then across populations, so that a large population cannot dominate
a small one and so that the cluster structure of Section~\ref{sec:metrics-inference}
is preserved.

\subsection{Point accuracy}

\begin{equation}
  \RMSE = \Big(\overline{(\hat y_{x,h}-y_{x,h})^2}\Big)^{1/2},\qquad
  \MAE = \overline{|\hat y_{x,h}-y_{x,h}|},\qquad
  \hat y_{x,h} = \hat F^{-1}_{x,h}(\tfrac12),
\end{equation}
together with the error in period life expectancy at birth,
$\hat e_0 - e_0$, computed from the median path. These are the metrics the
literature reports; \Hone\ is the claim that ranking by them is not ranking
by a proper score.

\subsection{Continuous ranked probability score}

For a predictive distribution $F$ and outcome $y$,
\begin{equation}
  \CRPS(F,y) = \int_{-\infty}^{\infty}\big(F(u)-\ind\{u\ge y\}\big)^2\,du
  = \E_F|Z-y| - \tfrac12\,\E_F|Z-Z'|,
\end{equation}
with $Z,Z'\sim F$ independent \citep{gneiting2007strictly}. From samples we
use the plug-in estimator with sorted samples $z_{(1)}\le\dots\le z_{(n)}$,
\begin{equation}
  \widehat{\CRPS} = \frac1n\sum_{j=1}^{n}|z^{(j)}-y|
  - \frac{1}{n^2}\sum_{j=1}^{n}\big(2j-n-1\big)\,z_{(j)},
  \label{eq:crps-sample}
\end{equation}
which is the $O(n\log n)$ form of the second term. CRPS is computed on the
log-rate scale, $y=y_{x,h}$, for every cell, and on the count scale as a
registered sensitivity (Section~\ref{sec:metrics-logscore}). It is strictly
proper and reduces to MAE for a point forecast, which is what makes the
RMSE-versus-CRPS ranking comparison of \Hone\ well posed.

\subsection{Poisson log score and the rounding convention}
\label{sec:metrics-logscore}

Each sample path implies a Poisson mean $\lambda^{(j)}_{x,h} =
E_{x,T+h}\,m^{(j)}_{x,T+h}$ for the death count, so the predictive count
distribution is the equal-weight mixture
$p(d)=\frac1n\sum_j \mathrm{Poisson}(d;\lambda^{(j)}_{x,h})$. The log score is
\begin{equation}
  \LS = -\log\Big(\frac1n\sum_{j=1}^{n}\mathrm{Poisson}\big(\tilde D_{x,T+h};\lambda^{(j)}_{x,h}\big)\Big),
  \qquad \tilde D = \big\lfloor D+\tfrac12\big\rfloor,
\end{equation}
evaluated by log-sum-exp. Because HMD deaths are fractional
(Section~\ref{sec:data}), the score is defined at the nearest integer;
the sensitivity is the CRPS of the mixture count distribution against the
unrounded $D$, which is well defined for any real $D$. Both conventions were
fixed before estimation.

\subsection{Interval score and coverage}

For the central $(1-\alpha)$ interval $[l,u] =
[\hat F^{-1}(\alpha/2),\hat F^{-1}(1-\alpha/2)]$ the interval (Winkler) score
is
\begin{equation}
  \IS_\alpha(l,u;y) = (u-l) + \frac{2}{\alpha}(l-y)\,\ind\{y<l\}
  + \frac{2}{\alpha}(y-u)\,\ind\{y>u\},
  \label{eq:winkler}
\end{equation}
\citep{winkler1972decision,gneiting2007strictly}, reported at $\alpha\in\{0.5,0.2,0.05\}$.
It is the width-penalised criterion whose absence from the mortality
literature Section~\ref{sec:intro} documents: an interval wide enough to
cover everything has coverage one and an interval score equal to its width.
Alongside it we always report the pair (empirical coverage, mean width),
never coverage alone.

\subsection{Marginal and joint path coverage}
\label{sec:metrics-coverage}

\emph{Marginal} coverage is the pointwise hit rate,
\begin{equation}
  C^{\mathrm{marg}}_{1-\alpha}(h) = \overline{\ind\{\,l_{x,h}\le y_{x,h}\le u_{x,h}\,\}},
\end{equation}
averaged over ages and populations for each horizon $h$, and pooled over
$h$ for a single number. It is the quantity every mortality paper that
reports coverage reports.

\emph{Joint path} coverage asks whether the whole trajectory of a unit
(one population--sex--age) lies inside its bands:
\begin{equation}
  C^{\mathrm{joint}}_{1-\alpha} = \overline{\prod_{h=1}^{H}\ind\{\,l_{x,h}\le y_{x,h}\le u_{x,h}\,\}},
  \qquad H=5 \text{ (shift, stable)},\ H=9 \text{ (placebo)}.
\end{equation}
A path counts as covered only if every horizon is inside its band. This is
the object an annuity valuation actually depends on, because a liability is
a functional of the path of future mortality, not of one year of it. For
pointwise bands at level $1-\alpha$, joint coverage is bounded above by
marginal coverage and, if horizons were independent, would equal
$(1-\alpha)^H$; the strong horizon dependence of random-walk paths puts it
between the two. Only the copula conformal mechanism targets the joint level
by construction. \Hthree\ is the registered claim that the gap
$C^{\mathrm{marg}}-C^{\mathrm{joint}}$ is material for every family, tested
as a paired difference over units.

\subsection{Randomised probability integral transform}

A forecaster is probabilistically calibrated if $F(Y)$ is uniform
\citep{gneiting2007probabilistic,dawid1984present}. Because the predictive is
an empirical distribution of $n$ samples and the target may tie with them,
we use the randomised PIT of \citet{czado2009predictive},
\begin{equation}
  u_{x,h} = \frac{\#\{j: z^{(j)}_{x,h} < y_{x,h}\} + V\,\big(\#\{j: z^{(j)}_{x,h}=y_{x,h}\}+1\big)}{n+1},
  \qquad V\sim\mathrm{Uniform}(0,1),
\end{equation}
which is exactly uniform under the null for any finite $n$. We report the
ten-bin histogram, the Kolmogorov--Smirnov distance from uniformity, and the
shape---U (over-confident), hump (under-confident), slope (biased). Because
PIT values within a population--sex--origin are strongly dependent, the
formal test is a cluster-bootstrap test of the pooled histogram
(Section~\ref{sec:metrics-inference}), not a cell-level KS $p$-value.

\subsection{Calibration by age}

All coverage and PIT statistics are additionally computed within the age
bands $0$--$24$, $25$--$64$ and $65$--$99$ and as single-age curves.
\Hfour\ compares the $65$--$99$ and $25$--$64$ bands in the shift regime.

\subsection{Murphy decomposition}

For a proper score $S$ the mean score decomposes as
\begin{equation}
  \overline{S} = \mathrm{UNC} - \mathrm{RES} + \mathrm{MCB},
\end{equation}
uncertainty of the outcome minus resolution of the forecasts plus
miscalibration \citep{murphy1973new,pohle2020murphy}. It separates
\emph{why} a forecaster scores badly: a well-calibrated but vague
forecaster has low RES; a sharp but wrong one has high MCB. We compute it
for the interval score and for the Brier scores of the exceedance events
$\ind\{Y\le \hat F^{-1}(p)\}$ at $p\in\{0.025,0.25,0.5,0.75,0.975\}$, with
the recalibrated forecast estimated by isotonic regression
\citep{dimitriadis2021stable}. \todo{estimator not yet implemented in the
harness; confirm the isotonic (CORP) choice or a binned alternative before
the first real-data run}

\subsection{Actuarial functionals}
\label{sec:metrics-actuarial}

Each sample path's rate vector $m^{(j)}_{\cdot,T+h}$ is converted to a period
life table by the conventions of Section~\ref{sec:data}, giving samples of
\begin{equation}
  e^{(j)}_x = \frac{T^{(j)}_x}{l^{(j)}_x},\qquad
  \annuity^{(j)} = \sum_{t\ge0} v^{t}\,{}_t p^{(j)}_{65},\quad v=\frac{1}{1.02},
\end{equation}
where ${}_t p_{65}=l_{65+t}/l_{65}$ from the sample's own table treated as
static, i.e.\ the annuity-due factor at 2\% on the period table of year
$T+h$ \citep{dickson2020actuarial}. The realised values are computed from the
observed $m_{x,T+h}$ by the same routine, so that any life-table
approximation cancels in the comparison. For $e_0$, $e_{65}$ and $\annuity$
we report the interval error, the empirical coverage of the central 95\%
interval and its width, in each regime; \Hfive\ compares the shift and
stable departures from nominal. We also report the 99.5\% upper quantile of
$\annuity$ from each cell as the longevity tail an internal model would
carry. \todo{the 99.5\% tail quantile is descriptive and not part of a
registered hypothesis; say so in the table note}

\subsection{Forecast-comparison inference}
\label{sec:metrics-inference}

\paragraph{Effective sample size.} The shift regime contains one common
shock. Within a population, ages and horizons share the same realised
pandemic; across populations, the shock itself is shared. Treating cells as
independent would produce standard errors that are wrong by orders of
magnitude. Population is therefore the cluster ($G=20$ in the shift and
stable regimes, $G=11$ in the placebo), and every test resamples clusters.
Even so, $p$-values in the shift regime are statements about differences
between forecasters \emph{given this event}, not about future pandemics;
that is what an event study can deliver, and the placebo regime is the
second event that makes a cross-event comparison possible.

\paragraph{Diebold--Mariano with a wild cluster bootstrap.} For forecasters
$A$ and $B$ with per-unit losses $L_{A,i}$, $L_{B,i}$ (CRPS per horizon per
population--sex--age), the differential $d_i = L_{A,i}-L_{B,i}$ has mean
$\bar d$ and cluster-robust variance
\begin{equation}
  \widehat{\mathrm{Var}}(\bar d) = \frac{G}{G-1}\,\frac{1}{N^2}\sum_{g=1}^{G}
  \Big(\sum_{i\in g}(d_i-\bar d)\Big)^2,
\end{equation}
giving the statistic $t=\bar d/\widehat{\mathrm{se}}$
\citep{diebold1995comparing}. Its null distribution is obtained by the wild
cluster bootstrap of \citet{cameron2008bootstrap} with the null imposed:
residuals $d_i-\bar d$ are multiplied by a cluster-level weight $w_g$ drawn
from the six-point distribution of \citet{webb2023reworking}, which has
enough distinct sign patterns at $G=20$ where Rademacher weights do not
\citep{mackinnon2017wild}; $B=4999$ replications. Negative $\bar d$ favours
$A$.

\paragraph{Model confidence set.} Following \citet{hansen2011model}, the MCS
at confidence $1-\alpha=0.90$ is built by the $T_{\max}$ statistic and the
$e_{\max}$ elimination rule, with the bootstrap distribution of the
statistic obtained by resampling \emph{populations} as blocks---the analogue
of their stationary block bootstrap for our panel. The MCS is reported for
CRPS on log rates in each regime, over the distributional cells of the grid.

\paragraph{Hypothesis tests.} \Hone\ is assessed by the Spearman correlation
of RMSE and CRPS rankings and the count of top-3 inversions, with a
cluster-bootstrap interval for the correlation. \Htwo, \Hfour\ and \Hfive\
are paired contrasts of coverage (shift minus stable; ages 65--99 minus
25--64) tested by the wild cluster bootstrap above applied to the coverage
indicators. \Hthree\ is the paired difference of marginal and joint
coverage over units, clustered by population.

\section{Results}
\label{sec:results}

\noindent\textit{Every table in this section is a generated fragment
(\texttt{paper/tables/<name>.tex}, one \texttt{booktabs} body per hook)
produced by the analysis stage from the per-cell results files; every
figure is a generated \texttt{figures/<name>.pdf}. No number in this
section is typed by hand. The shift-regime and placebo columns are final:
each is read from its regime's two results files
(\texttt{results/<regime>.parquet} and
\texttt{results/<regime>\_gp.parquet}) after that regime's post-sweep QA
gate passed. The stable columns of every regime-indexed table are pending
placeholders---that regime's results files do not yet exist in final
form---and the fragments carry a first-line marker to that effect until
all three regimes are present. Each table's note states the regime, the
number of populations, and---where a contrast is involved---the number of
cells dropped by the common-cell restriction of Addendum~3 (\S11).}

\subsection{Validation gates}

Gates 1--3 (Section~\ref{sec:design-gates}) passed before any real-data
result was read. The one substantive finding from gate~1---the mild
over-dispersion of two-stage SVD Lee--Carter on Poisson data---is carried
into the reading of the SVD-LC rows below. A fourth, post-sweep gate
(\texttt{scripts/final\_qa.py}) classifies every cell that produced no
valid row before any table is generated; its output is the first table of
this section.

\subsection{What was scored: design-floor and method-failure cells}
\label{sec:results-infeasible}

% tab-infeasible is a page-spanning longtable fragment generated by
% scripts/make_tables.py; its caption and label live in the fragment.
\inputtable{tab-infeasible}

Table~\ref{tab:infeasible} lists every (population, family, mechanism) cell
of the primary grid that produced no valid row, in the classes of the
QA gate; the unabridged ledger, one row per affected cell with the full
recorded reason, is Appendix~\ref{app:infeasible}. The shift sweep
produced 2{,}000 cell rows, of which 210 are error rows: 126 structural,
56 method, 28 in the gate's residual class, and none machine---the gate
that refuses to generate tables while a broken-run row exists passed on
the final files.

\emph{Structural} cells are design-floor cells: the panel is shorter than
the mechanism's calibration or ensemble-member requirement, or falls under
the registered admissibility floor of $n_{\text{train}}\ge 15$ complete
training years (Addendum~3, \S4). The 126 structural rows trace the
training lengths of Table~\ref{tab:populations} exactly: KOR, the shortest
shift panel, appears under every mechanism with a calibration demand; HRV
joins it wherever the demand grows from the split wrapper's single
calibration block to the ten refits of EnbPI and the copula wrapper; CHL
enters under the largest windowed demands (the EnbPI and copula arms of
the sparse VAR, CNN-LC and LSTM), and the LSTM's construction alone adds
HKG.
These cells are a property of the design, not a failure of any method.

\emph{Method} cells are those in which a family's own sampler declines to
produce a predictive law on that panel under a registered rule. In the
shift regime this class is exclusively the sparse VAR's
explosive-coefficient-draw rejection (Addendum~3, \S7---rejection and
never clipping): it fires on 11 rows over six populations under the
native law (CHL, HKG, HRV, JPN, KOR, TWN), on 15 rows over nine
populations under the Poisson bootstrap, and on 8--11 rows in each
conformal arm; these rows are read below as a property of the family
under its own registered rule. The Poisson-composition overflow guard
never fired.

The 28 residual rows (class \emph{other}) all belong to the multi-output
GP: cells whose trailing block of complete years is shorter than the
family's minimum training block (Section~\ref{sec:design-families}; the
sixty-year cap of Addendum~4 is Section~\ref{sec:design-deviations}).
They are design-floor in character---a family-specific admissibility
requirement rather than a sampler failure---and are classed apart because
the floor is the family's own, not the registered mechanism floor.

Error rows of every class are excluded from every mean in this paper, and
every contrast between arms below is computed on the intersection of cells
in which each compared arm produced a valid row, with the intersection
size reported (Addendum~3, \S11).

\subsection{Calibration where nothing breaks: the stable regime}

The stable regime has not yet been run: the stable rows of every
regime-indexed table in this section are a pending placeholder until
\texttt{results/stable.parquet} exists, and no stable-derived quantity---
including the second, comparative clause of \Htwo\ and \Hfive---is stated
as a result in this draft.

\subsection{H1 as a harness-validity check: ranking by RMSE versus ranking by proper scores}

% tab-h1-rankings is a page-spanning longtable fragment generated by
% scripts/make_tables.py; its caption and label live in the fragment.
\inputtable{tab-h1-rankings}

Table~\ref{tab:h1-rankings} ranks the distributional arms---the
family--mechanism cells whose samples constitute a predictive
distribution---three times within the shift regime: by mean RMSE on log
rates, by mean CRPS on log rates and by mean Poisson log score. Two
exclusions are by construction. The conformal arms do not appear, because
their CRPS, log score and PIT are placeholder scores from
uniform-in-interval samples and are never ranked against a distributional
mechanism. CBD is ranked in a separate block: it is fit on ages 55--99
and therefore scored on 45 ages, and a mean over a different age support
is not the same quantity as a mean over ages 0--99. Means are taken
within population--sex first and then across populations, never
exposure-weighted, and the ranking is computed on the common cells of the
block (22 of 40 units kept; the dropped 18 trace the sparse VAR, GP and
LSTM failure cells of Table~\ref{tab:infeasible}).

\Hone\ as registered---rankings by RMSE and by CRPS disagree, with rank
correlation below one and at least one top-three inversion, in the shift
regime---is supported, in its mild form against CRPS and in a much
stronger form against the log score. The Spearman correlation between the
RMSE and CRPS orderings is $0.95$: below one, with the registered
top-three inversion present---the RMSE top three (LSTM-$\kappa_t$ under
dropout and ensemble, then SVD-LC native) share only one member with the
CRPS top three (the two sparse VAR arms, then LSTM dropout). Against the
Poisson log score the point-accuracy ordering nearly dissolves:
$\rho=0.12$. The arms with the best RMSE are among the worst by log
score---the LSTM ensemble is 2nd by RMSE and 16th of 18 by log score,
SVD-LC native 3rd and 12th---while the NB head, mid-pack by RMSE (ranks
10--12), takes the top three log-score positions together with CNN-LC
dropout, and neural-LC dropout, last by RMSE, is 5th by log score. A
point-accuracy league table and a density-forecast league table of the
same arms over the same break are close to unrelated, which is the
inversion the published literature reports
\citep{barigou2023bayesian,goes2025bayesian,schnurch2022point} recurring
under pre-registration.
% Intervals: results/h5_extras.json (scripts/h5_extras.py) -- cluster
% bootstrap resampling populations with replacement on the tab-h1
% common-cell block, B = 2000, seed 20260830.
Both correlations carry the cluster-bootstrap interval of
Section~\ref{sec:metrics-inference}, obtained by resampling with
replacement the thirteen populations that survive the common-cell
restriction and recomputing the block's means, ranks and correlations in
each of 2{,}000 replicates: $[0.77, 0.96]$ for $\rho(\text{RMSE},
\text{CRPS})$ and $[-0.16, 0.21]$ for $\rho(\text{RMSE}, \text{log
score})$. The first interval excludes one, so the registered
below-one clause does not depend on which populations entered the
block; the second contains zero, so at the resolution thirteen
correlated clusters afford, the log-score ordering is indistinguishable
from one unrelated to the RMSE ordering.

\subsection{H2: marginal coverage across the break}

\begin{table}[p]\centering
\caption{Empirical coverage of nominal 50/80/95\% central intervals and
mean 95\% interval score, by family, mechanism and regime, with the number
of cells behind each mean. Conformal mechanisms are scored at their
construction level only, so their 50\% and 80\% columns are blank by
design (Addendum~2, \S3). Error rows are excluded from every mean and
counted in Table~\ref{tab:infeasible}; a regime whose results file does
not yet exist appears as a pending row.}
\label{tab:h2-coverage}
\inputtable{tab-h2-coverage}
\end{table}

Table~\ref{tab:h2-coverage} reports, for every admissible family--mechanism
cell, the mean empirical coverage at 50, 80 and 95\%, the mean 95\%
interval score, and the number of cells behind each mean; conformal rows
carry the 95\% columns only, because a conformal wrapper constructs one
interval at one level and its coverage and interval score are read from
that interval's bounds. The interval score is the width-penalised
companion that keeps a vague interval from passing as a calibrated one.

\Htwo\ as registered has two clauses. The comparative clause---the
shortfall exceeds the stable-regime shortfall---awaits the stable regime
and is not tested here. The universal clause---nominal 95\% intervals
under-cover in the shift regime \emph{for every model family}---is
refuted as a universal, and the refutation is the more informative
result. The classical cores under their own predictive law do under-cover,
severely: Poisson Lee--Carter 0.624, APC/RH 0.654, SVD-LC 0.692, CBD
0.799 (on its 45-age support). But the distributional NB head covers
0.933 natively, 0.956 under the deep ensemble and 0.972 under dropout;
CNN-LC under dropout covers 0.962; the multi-output GP 0.905; the sparse
VAR 0.937--0.938 on the cells its rejection rule admits. Three of those
arms sit at or above nominal, so the direction of the miscalibration is
not even uniformly downward.

The registration deliberately declined a directional claim on
neural-versus-classical ordering, and the working conjecture that
motivated the audit---coverage degradation larger for the neural
families---is refuted in its strong form. The split does not run along
the architecture class at all: LSTM-$\kappa_t$ (0.647 ensemble, 0.685
dropout) and the neural-LC ensemble (0.694) break exactly like the
classical cores, while their sibling neural arms with a count likelihood
or dropout smoothing (NB head, CNN-LC dropout) are the best-covered arms
in the study. What separates the broken arms from the calibrated ones is
the family--mechanism pair, not the presence of a neural network---the
finding the crossed design exists to expose, taken up in
Section~\ref{sec:results-contrasts}.

The interval score orders the same failure: the broken natives pay twice,
in missed observations and in width---Poisson Lee--Carter native 4.456,
neural-LC ensemble 4.321, APC/RH 3.401, against NB native 1.684 and
CNN-LC dropout 1.677. The largest interval score in the table, neural-LC
dropout at 5.559 with coverage 0.811, is an arm buying part of its
coverage with indiscriminate width. Wrapping any family in a conformal
band moves its 95\% coverage into a narrow range regardless of how broken
the native law is: 0.848--0.944 for split and EnbPI, 0.939--0.981 for
the copula arm---the repair whose cost is priced in
Table~\ref{tab:dm-mcs}.

\subsection{H3: joint path coverage}

\begin{table}[p]\centering
\caption{Marginal coverage of nominal 95\% intervals, pooled over horizons,
against joint path coverage over the whole horizon set ($H=5$ in the shift
and stable regimes, $H=9$ in the placebo), and their difference, by family,
mechanism and regime. Conformal rows enter with both quantities read from
their interval bounds at the construction level (Addendum~3, \S6).}
\label{tab:h3-joint}
\inputtable{tab-h3-joint}
\end{table}

Table~\ref{tab:h3-joint} sets, for each family--mechanism cell, mean
marginal coverage at 95\% beside mean joint path coverage at the same
nominal level and the gap between the two: the difference between the
coverage a mortality paper reports and the coverage a liability depends
on. The registered direction of \Hthree---joint path coverage materially
below marginal for every family---holds in every row of the table: the
gap is negative for all 50 admissible arms, without exception. For the
broken classical natives the collapse is near-total: Poisson Lee--Carter
falls from 0.624 marginal to 0.298 joint, SVD-LC from 0.692 to 0.417,
APC/RH from 0.654 to 0.330, CBD from 0.799 to 0.504. The gap is not a
symptom of bad marginal calibration that good calibration removes: the
best-covered marginal arms lose it too---NB dropout falls from 0.972 to
0.897, CNN-LC dropout from 0.962 to 0.870---so an arm whose reported
marginal coverage is nominal can still miss between one in ten and one in
three whole trajectories.

The one mechanism constructed to be joint over $h=1,\dots,H$ behaves as
constructed: the copula-conformal arm holds joint path coverage between
0.805 (CBD) and 0.927 (SVD-LC) across the ten families, with gaps of
$-0.049$ to $-0.136$---the smallest in its family block in every family
but the NB head, whose dropout arm's gap of $-0.075$ is narrower still.
The registered comparator for \Hthree\ is not the marginal rate but each
model's own simulated-path-implied joint coverage (Addendum~3, \S8),
computed from the same predictive paths, so that the
realised-minus-implied gap can take either sign; that quantity is not
among the columns the runner emits. \todo{model-implied joint coverage
column (Addendum~3 \S8) is not emitted by the runner; add it to the
analysis stage or record the omission as a deviation in
Section~\ref{sec:design-deviations}}

\begin{figure}[htbp]\centering
\inputfigure{fig-joint-vs-marginal-shift}
\caption{Marginal 95\% coverage against joint path coverage over
$h=1,\dots,H$ as paired bars per mechanism, one panel per family, shift
regime, on the common cells of each panel (intersection size in the panel
title). Conformal arms are drawn in a distinct style and read from their
interval bounds; the nominal level is marked. The stable and placebo
counterparts are the same figure on those regimes' files.}
\label{fig:joint-vs-marginal}
\end{figure}

\subsection{H4: where in age the intervals fail}

\begin{table}[p]\centering
\caption{Empirical coverage of nominal 95\% intervals within the age bands
$0$--$24$, $25$--$64$ and $65$--$99$, by family, mechanism and regime.
CBD is fit on ages 55--99: its $0$--$24$ entry is blank and its
$25$--$64$ entry covers ages 55--64 only, as the table note states.
Conformal rows are read at the same bands the wrappers calibrate on.}
\label{tab:h4-age}
\inputtable{tab-h4-age}
\end{table}

Table~\ref{tab:h4-age} gives, for each family--mechanism cell, the 95\%
coverage within the three registered age bands, so that the $65$--$99$
and $25$--$64$ columns can be read against each other as \Hfour\
requires; the conformal wrappers are Mondrian-stratified on the same
bands, so a conformal row is read at the resolution it was calibrated at.
CBD contributes no $0$--$24$ value and a $25$--$64$ value restricted to
ages 55--64, which the table note records so that its row is not read
against the full-age families.

\Hfour\ as registered---shift-regime coverage at ages 65--99 lower than
at 25--64 \emph{for every family}---is refuted as a universal: the age
profile of the failure is family-specific, and no single clean pattern
spans the table. The registered direction does hold, sharply, for the
three families that forecast through an extrapolated period index and
broke in Table~\ref{tab:h2-coverage}: SVD-LC native falls from 0.725 at
25--64 to 0.518 at 65--99, APC/RH from 0.661 to 0.589, and
LSTM-$\kappa_t$ most steeply of all, 0.687 to 0.426 (ensemble) and 0.733
to 0.463 (dropout), against markedly better 0--24 rows in each case
(0.882--0.920 for SVD-LC and the LSTM, 0.732--0.744 for APC/RH): for
these arms the COVID miss is concentrated where COVID mortality was, at
the old ages. But Poisson Lee--Carter is uniformly
poor across all three bands (0.639 / 0.602 / 0.639)---its failure has no
age profile---and three arms reverse the registered direction outright:
CBD covers 0.849 at 65--99 against 0.626 at its 55--64 stub, CNN-LC
ensemble 0.880 against 0.847, and neural-LC dropout 0.934 against 0.874.
The neural-LC dropout arm instead collapses at the young ages: 0.538 at
0--24, the lowest band entry of any arm outside the index-driven
families. The calibrated families (NB head, sparse VAR, GP) are
near-flat in age, and the Mondrian-stratified wrappers flatten even the
broken families' profiles (SVD-LC split: 0.944 / 0.897 / 0.945). The
reversal anticipated by \citet{dowd2010backtesting}, which Addendum~3
(\S8) pre-committed to reporting as informative, is therefore present in
the table---but only for particular family--mechanism pairs, not as the
complement of a uniform old-age failure.

\begin{figure}[htbp]\centering
\inputfigure{fig-coverage-by-age-shift}
\caption{Single-age empirical coverage of nominal 95\% intervals (mean
over horizons), one panel per family, mechanisms as lines with conformal
arms dashed, shift regime, on the common cells of each panel. Ages a
family does not score are left blank and CBD's panel states its age
support; the horizontal line is the nominal level. The stable and placebo
counterparts are the same figure on those regimes' files.}
\label{fig:coverage-by-age}
\end{figure}

\subsection{PIT uniformity, reliability and the Murphy decomposition}

\begin{table}[htbp]\centering
\caption{Distributional arms only: mean Kolmogorov--Smirnov distance of the
randomised PIT from uniformity, and the share of population--sex cells
whose nominal KS $p$-value falls below $0.05$, by family, mechanism and
regime; conformal rows are left blank because their PIT is a placeholder.
The $p$-value is descriptive: PIT values within a cell are dependent
across ages and horizons (Addendum~2, \S4), and formal inference on
calibration is population-clustered.}
\label{tab:pit}
\inputtable{tab-pit}
\end{table}

\begin{table}[p]\centering
\caption{Murphy decomposition of the 95\% interval-hit Brier score
(reliability, resolution, uncertainty) and the PIT-scale reliability
against the uniform reference, by family, mechanism and regime; with one
constant nominal level the resolution term is zero by construction
(table footnote). PIT-scale column: distributional arms only.}
\label{tab:murphy}
\inputtable{tab-murphy}
\end{table}

Table~\ref{tab:pit} reports, for the distributional arms, the mean KS
distance of the randomised PIT from the uniform and the share of cells
whose nominal KS $p$-value is below $0.05$; the $p$-value is descriptive
only, because PIT values across ages and horizons within a cell are
dependent, and the shape of each departure is read from
Figure~\ref{fig:pit-hist}. The ordering matches
Table~\ref{tab:h2-coverage}---the smallest KS distances belong to the GP
(0.172), the NB head (0.182--0.189) and the sparse VAR (0.188), the
largest to Poisson Lee--Carter (0.374) and APC/RH (0.325)---but the
levels carry the sharper message: no arm is distributionally calibrated
through the break. Even for the best arms the share of cells whose
nominal KS $p$-value falls below 0.05 never drops beneath 0.844, and for
most arms it is 0.925--1.000. Coverage at a single level is a much
weaker property than PIT uniformity: the NB head's 95\% interval can be
essentially nominal while its predictive distribution, read over all
levels at once, is still rejected as uniform in nearly every cell.

Table~\ref{tab:murphy} decomposes the Brier score of the 95\% hit
indicator: with a constant forecast probability there is one bin, so
reliability is the squared coverage gap, resolution is zero by
construction, and uncertainty is the base-rate term. The reliability
column spans two orders of magnitude across arms---0.1515 for Poisson
Lee--Carter native, 0.1311 for APC/RH and 0.1226 for the LSTM ensemble
at the broken end, against 0.0015 for CNN-LC dropout, 0.0022 for NB
dropout and 0.0012--0.0223 for the conformal rows---which restates the
coverage table as a proper-score share. The PIT-scale reliability beside
it, the chi-square-type distance from uniformity that a single level's
coverage cannot see, orders the arms the same way (Poisson Lee--Carter
0.1904 down to NB dropout 0.0302) but never approaches zero, in
agreement with the KS column of Table~\ref{tab:pit}: the calibrated arms
are calibrated at 95\%, not everywhere. \todo{Section~\ref{sec:metrics}
describes a Murphy decomposition of the interval score and of exceedance
Brier scores with an isotonic recalibration; the harness emits the
one-bin hit-indicator form and the PIT-scale form above---reconcile the
metrics text with the emitted columns before submission}

\begin{figure}[htbp]\centering
\inputfigure{fig-pit-hist-shift}
\caption{Ten-bin randomised PIT histograms averaged across population--sex
cells, in a family $\times$ distributional-mechanism grid, shift regime;
the horizontal line is the uniform reference. Conformal arms are absent
because their PIT is a placeholder. A U shape reads as over-confidence, a
hump as under-confidence, a slope as bias. The stable and placebo
counterparts are the same figure on those regimes' files.}
\label{fig:pit-hist}
\end{figure}

\begin{figure}[htbp]\centering
\inputfigure{fig-reliability-shift}
\caption{Reliability diagram: empirical coverage against nominal level
(50, 80, 95\%) per mechanism, one panel per family, shift regime, with
the diagonal as perfect calibration. Conformal arms contribute a single
point at their construction level, 95\%, in a distinct style. The stable
and placebo counterparts are the same figure on those regimes' files.}
\label{fig:reliability}
\end{figure}

\subsection{Forecast comparison: Diebold--Mariano and the model confidence set}

% tab-dm-mcs is a page-spanning longtable fragment generated by
% scripts/make_tables.py; it carries its own \caption and \label{tab:dm-mcs}.
\inputtable{tab-dm-mcs}

Table~\ref{tab:dm-mcs} shows the two registered inference procedures, and
its loss column is the reason it is one table: a contrast among
distributional arms is decided on CRPS on log rates, while any contrast
that includes a conformal arm is decided on the per-horizon 95\% interval
score, because a conformal cell's CRPS is a placeholder. Every entry is
computed on the intersection of cells in which each compared arm produced
a valid row, and the kept and dropped counts are part of the table.

The model confidence set over the classical families under their own
predictive law is formed over the four full-age families only: the
five-family set including CBD is skipped, with its reason recorded in the
table, because a mean CRPS over 45 ages is not comparable with one over
100 and the registration forbids working around an incomparable loss. On
the 29 common cells the 90\% set retains \{SVD-LC native, sparse VAR
native\} (elimination $p$-values 0.13 and 1.00), eliminating Poisson
Lee--Carter and APC/RH at $p=0.00$---with the caveat, carried from
Table~\ref{tab:infeasible}, that the sparse VAR is scored only on the
cells its own rejection rule admits.

The within-family conformal sets on the interval score split three ways.
In five families---SVD-LC, LSTM-$\kappa_t$, NB head, neural-LC and sparse
VAR---EnbPI is the sole survivor, eliminating the split and copula arms
at $p$-values between 0.00 and 0.08; in Poisson Lee--Carter the copula
arm is the sole survivor ($p=0.01$ against both others); and in CBD,
CNN-LC, the GP and APC/RH all three wrappers are retained. Read with
Table~\ref{tab:h3-joint}, the EnbPI wins are a width result, not a
calibration result: the interval score at a single level rewards the
narrowest band that roughly covers, which EnbPI's residual recycling
produces, while the copula arm's wider joint band is what actually holds
path coverage. The two registered procedures answer different questions
and disagree accordingly.

The Diebold--Mariano block prices the conformal repair of
Table~\ref{tab:h2-coverage} and is the attribution finding of the crossed
design. The contrast is defined for the seven families with a native
arm. Split conformal is significantly \emph{better} than the native law
for exactly the four classical cores whose native intervals broke:
Poisson Lee--Carter ($\Delta=1.895$, $p=0.004$), APC/RH ($\Delta=1.039$,
$p=0.013$), SVD-LC ($\Delta=0.700$, $p=0.018$) and CBD ($\Delta=0.318$,
$p=0.017$), with $G=19$ population clusters each. It is significantly
\emph{worse} for exactly the three families whose native law was already
calibrated: the NB head ($\Delta=-1.024$, $p=0.000$, $G=19$), the
multi-output GP ($\Delta=-0.413$, $p=0.009$, $G=16$) and the sparse VAR
($\Delta=-0.318$, $p=0.000$, $G=14$). No family is left where the
wrapper is free: conformalisation is insurance that pays out precisely
when the model's own uncertainty was wrong and costs a premium precisely
when it was right. Which side of the line a family falls on is a
property of the family--mechanism pair, not of the wrapper---the
sharpest form of the study's central claim that architecture and
uncertainty mechanism are separate factors.

\subsection{Twin crises: 1914--1922 beside 2020--2024}

\begin{table}[htbp]\centering
\caption{Placebo (train to 1913, test 1914--1922, eleven populations) and
shift (train to 2019, test 2020--2024, twenty populations) side by side:
marginal 95\% and joint path coverage by family and mechanism, with the
placebo strata of Addendum~1 as a further column. Joint coverage spans
$H=9$ placebo and $H=5$ shift horizons---not one scale. Sparse VAR's
empty placebo cells are a finding, not a gap: every native draw on the
pre-1914 panels tripped the registered explosive-path rejection
(Table~\ref{tab:infeasible}).}
\label{tab:twin-crises}
\inputtable{tab-twin-crises}
\end{table}

Table~\ref{tab:twin-crises} sets the 1914--1922 placebo beside the
2020--2024 shift. Read within regime (the joint columns span $H=9$
against $H=5$), the placebo reproduces the shift regime's family split a
century earlier: the full-age classical cores under-cover (Lee--Carter
0.787, Poisson Lee--Carter 0.780, APC/RH 0.761 marginal against a
nominal 0.95), the point-forecast neural mechanisms are worst
(neural-LC ensemble 0.661, LSTM 0.636), and the families that came
through COVID calibrated sit near the top of the placebo ordering too
(NB head 0.849, multi-output GP 0.820). Joint path coverage collapses
in the same order---0.146 (neural-LC ensemble), 0.172 (LSTM),
0.293--0.335 (the classical cores) over nine horizons---so the H3 gap
is not a COVID artefact either. Sparse VAR is the one family with no
placebo verdict: every native draw on the pre-1914 panels tripped the
registered explosive-path rejection (zero valid rows,
Table~\ref{tab:infeasible}), which also empties the common-cell
intersection (addendum 3~\S11) of the registered classical model
confidence set in this regime.

The exception is CBD, and it is the most instructive cell in the table:
near-nominal through the placebo break (0.986 marginal, 0.918 joint)
and broken through the COVID one (0.799, 0.504). The 1918 influenza
concentrated its excess mortality at young-adult ages, below CBD's
age-55 support, while COVID's excess fell at the old ages CBD models.
Whether an interval survives a crisis therefore depends not only on the
family's dynamics but on where in the age range the crisis lands---the
age-localisation of miscalibration that H4 establishes \emph{within}
the shift regime, operating here \emph{between} regimes.

The Addendum~1 strata sharpen the reading. In the neutral
register-based populations the classical cores hold 0.86--0.89---the
1918 flu alone bends them without breaking them---while the
belligerent-total stratum, where war losses compound the flu, drops the
same families to 0.55--0.59, with the civilian-only stratum in between.
Within one family and one regime, coverage failure scales with the size
of the break.

The paired mechanism contrasts, however, do not transfer. The
split-conformal rescue that is decisive under COVID
(Section~\ref{sec:results-contrasts}) is absent here: on the interval
score, native versus split is insignificant for every full-age
classical core (Lee--Carter $\Delta=-0.091$, $p=0.689$; Poisson
Lee--Carter $p=0.127$; APC/RH $p=0.563$), and the NB head's native law
significantly beats its own wrapper ($\Delta=-0.461$, $p=0.008$)---a
calibration window drawn from the pre-1914 years knows as little about
the coming shock as the model it corrects. The ensemble-versus-dropout
verdicts flip family by family between the two crises: CNN-LC,
indistinguishable under COVID ($p=0.739$), favours dropout here
($\Delta=+0.531$, $p=0.0002$); LSTM, which favours dropout under COVID
($p=0.0006$), is indistinguishable here ($p=0.478$). A family's
fragility is predictable from its past crisis; which uncertainty
mechanism patches it best is not.

\subsection{Architecture or mechanism? Contrasts within sub-grids}
\label{sec:results-contrasts}

The payoff of the crossed design. Each contrast below is pre-registered and
computed as a paired difference within the admissible sub-grid named, on
the common cells of that sub-grid.

\begin{enumerate}
\item \emph{Conformal versus native.} Does wrapping a fit in a conformal
  band change its shift-regime interval score and coverage? The native
  versus split-conformal block of Table~\ref{tab:dm-mcs} carries this
  contrast, per family, on the interval score, and its answer is signed by
  the family: the wrapper significantly improves the four classical cores
  whose native law broke and significantly worsens the NB head, GP and
  sparse VAR, whose native law had not. The width side of the question is
  read from the interval-score column of Table~\ref{tab:h2-coverage}: for
  the broken families the wrapper buys its coverage while \emph{shrinking}
  the score (Poisson Lee--Carter 4.456 native against 2.674 split), for
  the calibrated ones it pays in width (NB head 1.684 native against
  2.717 split).
\item \emph{Ensemble versus Monte Carlo dropout, neural families only.} Do
  two mechanisms on the same architecture give the same coverage?
  Descriptively, Table~\ref{tab:h2-coverage} answers no---on the same
  CNN-LC architecture, dropout covers 0.962 against the ensemble's 0.882,
  and on neural-LC 0.811 against 0.694. The registered paired test agrees
  and adds a sign that the descriptive reading does not supply: on the
  interval score the two mechanisms differ significantly on two of the four
  neural families, \emph{in opposite directions}. Dropout is favoured on
  LSTM-on-$\kappa_t$ ($\Delta = +0.224$, $t = 3.78$, $p = 0.0006$) and the
  ensemble on neural-LC ($\Delta = -1.238$, $t = -4.18$, $p < 0.0001$),
  while CNN-LC ($p = 0.739$) and the NB head ($p = 0.139$) are
  indistinguishable. The mechanism is therefore not separable from the
  architecture it wraps: no single ordering of the two survives a change of
  family, which is the crossed design's reason for existing.
\item \emph{Bootstrap versus native, classical families only.} Does adding
  parameter uncertainty move a classical interval toward nominal at the
  break? Descriptively the movement is positive and small (SVD-LC 0.692
  to 0.727, CBD 0.799 to 0.830 in Table~\ref{tab:h2-coverage}), nowhere
  near the repair a conformal wrapper delivers. The registered paired test
  makes the smallness precise: the semiparametric bootstrap improves the
  interval score significantly on three of the five classical families---
  SVD-LC ($\Delta = -0.112$, $p < 0.0001$), Poisson Lee--Carter
  ($\Delta = -0.115$, $p < 0.0001$) and Renshaw--Haberman
  ($\Delta = -0.101$, $p = 0.0002$)---leaves CBD unchanged ($p = 0.681$),
  and significantly \emph{worsens} sparse VAR ($\Delta = +0.168$,
  $p = 0.0008$, on the 22 cells where its native arm survives the
  explosive-draw rejection of Addendum~3~\S7). The three improvements are
  an order of magnitude smaller than the conformal repair of the same
  families in the contrast above (for Poisson Lee--Carter, $-0.115$ against
  $+1.895$ in favour of the wrapper): resampling the estimator moves a
  broken interval, but it does not fix it.
\item \emph{Joint-path conformal versus split conformal, all families.} The
  \Hthree\ mechanism contrast: the within-family conformal confidence sets
  of Table~\ref{tab:dm-mcs} on the interval score, read with the joint
  column of Table~\ref{tab:h3-joint}. The two tables pull in opposite
  directions by construction: the copula arm holds joint path coverage of
  0.805--0.927 where the split arm holds 0.692--0.862, yet on the
  single-level interval score the copula arm survives as sole set member
  only for Poisson Lee--Carter and is eliminated in five families where
  EnbPI's narrower band wins. A practitioner who needs the path pays for
  it at the level; the study reports the price rather than adjudicating
  it.
\end{enumerate}

\section{Actuarial impact}
\label{sec:actuarial}

\noindent\textit{Skeleton; all tables are placeholders.}

\subsection{H5: propagation into life expectancy and annuity factors}

\begin{table}[htbp]\centering
\caption{Coverage of nominal 95\% intervals, interval width and error for
$e_0$, $e_{65}$ and $\annuity$ at 2\%, by family and mechanism, stable versus
shift regime.}
\label{tab:h5-actuarial}
\inputtable{tab-h5-actuarial}
\end{table}

The table will show whether an interval failure on the rate surface
survives aggregation into the quantities that are actually valued---or
whether the life-table integration averages it away. Both outcomes are
informative: the first says the reserve is wrong, the second says the
audit's headline is a statement about rates rather than about liabilities.
\todo{result pending}

\subsection{Reading the failure in balance-sheet terms}

For an annuity-due of one per annum at age 65, the interval on $\annuity$
translates directly into an interval on the best-estimate liability of a
unit annuitant. A nominal 95\% interval that covers the realised
$\annuity$ in materially fewer than 95\% of population--years is, in
Solvency~II language, an internal model whose longevity distribution is
too narrow; the size of the shortfall in the 97.5\% quantile of $\annuity$
is the amount by which the implied capital at that level is understated.
We report this quantity as a proportion of the central $\annuity$ so that
it is portable across populations and interest-rate assumptions.
\todo{decide whether to present a worked example on a hypothetical
portfolio; if so, no figures until the H5 table exists}

\subsection{The longevity tail}

The 99.5\% upper quantile of the predictive distribution of $\annuity$ is
the one-year longevity tail an internal model would carry. We report it
alongside the realised $\annuity$ for each cell; the fraction of cells in
which the realised value exceeds the model's own 99.5\% quantile is the
empirical frequency of a nominally once-in-200-year event, and is the
quantity the pandemic years let us observe for the first time.
\todo{descriptive; not a registered hypothesis, label as such}

\subsection{Twin crises in economic terms}

The placebo regime supplies the same functionals for 1914--1922. A
side-by-side of the $e_{65}$ and $\annuity$ interval shortfalls across the
two events asks whether the economic size of a break-induced interval
failure is stable across a century, or whether the pandemic years are
distinctive. \todo{result pending; gated on the placebo data-quality check}

\section{Discussion}
\label{sec:discussion}

The limitations below are pre-committed: they were written before any
result was read, so that none can be selected after the fact to excuse a
finding or to inflate one.

\subsection{The effective sample size is one}

Every shift-regime number in this paper is a statement about one pandemic
observed in twenty places. Clustering by population and imposing the null
in the bootstrap make the standard errors honest \emph{given the event}; they
do not make the event a sample of events. The audit therefore answers
``did these intervals survive this break?'' and, through the placebo,
``does the pattern resemble the previous break?''---not ``will these
intervals survive the next break?''. A reader who wants the last question
answered wants a different study, one that we cannot run until there is a
next break.

\subsection{The grid is ragged, so claims are contrasts}

Ten families by seven mechanisms is seventy cells on paper and fifty in
fact. Main effects over a ragged grid confound the factor with the pattern
of admissibility, so no main effect is reported; every claim is a contrast
within a sub-grid whose composition is stated. This costs generality: a
finding that conformal wrapping changes coverage for classical families
says nothing about neural families unless the neural sub-grid shows it
too, and we will not average across the two.

\subsection{Conformal cells carry secondary proper scores}

A conformal band is an interval, not a distribution. Drawing uniformly
inside it keeps the single interface honest for coverage, width and
interval score, but the CRPS, log score and PIT of those draws describe a
placeholder density. They are reported in Appendix~\ref{app:secondary},
flagged, and not ranked. Any reader who compares a conformal cell's CRPS
with a native cell's CRPS is comparing a convention with a forecast.
Conformal predictive systems that output a full calibrated distribution
would remove this limitation and are noted as future work.

\subsection{2024 may be revised}

HMD does not flag provisional years; it re-issues whole series at each
release, so the 2024 values are final only in the sense that this release
contains them. The shift regime's $h=5$ results rest on 2024, and two of the
twenty populations (USA, CHL) carry census-based exposures that HMD
documents as pending revision. The pre-registered robustness run drops 2024;
Addendum~1 adds runs that drop USA and CHL and that contrast register-based
with census-based populations. If the $h=5$ and $h\le4$ results disagree,
the $h\le4$ result is the one to trust and the $h=5$ result is reported as
provisional; if a later vintage appears before submission, the shift regime
is re-run on it and any change in the headline coverage is reported as a
vintage-revision effect.

\subsection{SVD Lee--Carter is over-dispersed by construction}

Validation gate~1 found that the classical two-stage SVD Lee--Carter
produces intervals slightly wider than nominal on data simulated from a
Poisson Lee--Carter process, because the SVD estimate of $\kappa_t$ absorbs
observation noise that then inflates the random-walk innovation variance.
This is a property of the method that \citet{girosi2007understanding}
would predict, not a harness artefact, and it has a direct consequence for
reading the results: an SVD-LC row that appears better calibrated than
Poisson-LC in the stable regime may be so partly because it is vaguer, and
the interval score and Murphy resolution term---not coverage---are the
columns that tell the two apart.

\subsection{Other limitations}

\begin{itemize}
\item \emph{Exchangeability is violated by design.} Conformal calibration
  residuals come from pre-cutoff years, and the test years are precisely
  the years chosen to differ from them. \citet{barber2023conformal} bound the
  resulting coverage gap; we measure it rather than assume it away. The
  conformal arms are audited, not endorsed.
\item \emph{Neural families are budget-limited.} Ensembles of ten over
  thirteen origins, twenty populations and two sexes fix the hyperparameter
  search to what fits the GPU budget; the inner time split is small (five
  years). A better-tuned network might do better on point accuracy; whether
  it would be better calibrated is exactly what is not known.
\item \emph{Data quality in the placebo window.} Military deaths are
  reconstructed in the belligerent series and excluded in the civilian-only
  one. Placebo results are descriptive, are reported pooled and by the
  pre-specified strata, and the neutral register-based stratum carries the
  1918-versus-2020 comparison.
\item \emph{Single-age granularity.} Modelling ages $0$--$99$ singly keeps
  the old-age cells noisy in small populations; the age-cap-90 robustness
  run tests whether \Hfour\ is an artefact of that noise.
\item \emph{Rounding of deaths.} The log-score convention rounds fractional
  deaths; the CRPS-on-counts sensitivity exists to show that nothing turns
  on it.
\item \emph{One code path, one author.} A single evaluation module is a
  strength for consistency and a weakness for independent verification; the
  synthetic-truth gate and the \pkg{StMoMo} oracle are the checks, and the
  code is public.
\end{itemize}

\subsection{What this paper does not do}

It does not propose a model, a repair, or a recommendation about which
family to use. If the audit finds that a mechanism restores coverage at the
break, that is a finding about the mechanism to be stated as a contrast; it
is not a method contribution and will not be presented as one. If the audit
finds that nothing holds nominal coverage, that is the finding, and the
actuarial reading in Section~\ref{sec:actuarial} is its consequence.
Explainability of the neural families---which features drive their
forecasts and whether the explanations are stable across the break---is a
separate question deferred to a companion paper that reuses these fits.

\section{Conclusion}
\label{sec:conclusion}

Mortality prediction intervals are priced, reserved against and hedged,
and until now none had been scored against a realised structural break.
This paper pre-registered a calibration audit---ten forecasting families
crossed with seven uncertainty mechanisms, fifty admissible cells, one
sampling interface, one scoring path---and has now run its primary regime:
every interval built strictly from data through 2019, held against
realised 2020--2024 mortality in twenty populations of the first HMD
vintage with final data through 2024.

On the shift regime the registered hypotheses read as follows. \Hone,
demoted in advance to a harness-validity check, is supported: the RMSE and
CRPS orderings disagree with the registered top-three inversion, and
against the Poisson log score the point-accuracy ordering nearly dissolves
($\rho=0.12$). The universal clause of \Htwo\ is refuted in strong form:
the classical cores under their own predictive law covered 62--80\% at a
nominal 95\% (Poisson Lee--Carter 0.624, Renshaw--Haberman 0.654, SVD
Lee--Carter 0.692, CBD 0.799), but the NB distributional head
(0.933--0.972), the dropout CNN (0.962) and the sixty-year-window GP
(0.905) held near-nominal coverage through the break while the LSTM
(0.647--0.685) and the ensembled neural-LC (0.694) broke alongside the
classical cores; the failure is a property of the family--mechanism pair,
not of an architecture class. The registered direction of \Hthree\ holds
in every one of the fifty admissible arms: joint path coverage collapses
relative to marginal (0.298 against 0.624 for Poisson Lee--Carter), and
only the copula band, the one mechanism calibrated for the path, holds
0.805--0.927 joint. \Hfour\ is refuted as a universal: the old-age
concentration it registered is real but belongs to the index-extrapolating
families, while other arms are age-flat or reverse it. The shift half of
\Hfive\ is established---coverage of the registered $e_{65}$ and
$\annuity$ intervals departs from nominal by at least 0.200 for every
full-age distributional arm, and integration concentrates the rate-surface
failure rather than averaging it away. The comparative clauses of \Htwo\
and \Hfive\ await the stable control, and the 1914--1922 placebo supplies
the second event; both run under the same registered design and no claim
here rests on them.

The crossed design delivered the attribution it was built for: split
conformal wrapping significantly improves the interval score of exactly
the four classical cores whose native law broke and significantly worsens
exactly the three families whose native law had held. That sign pattern,
not any single coverage number, is the study's central result, and it is
the kind of result only a design that separates architecture from
uncertainty mechanism can produce.

The remaining contribution is the protocol itself: a pre-registration
hash-stamped before the first fit, a harness that recovered nominal
coverage on synthetic truth and matched the reference implementation to
machine precision, one predictive-sample interface for every family, and
cluster-bootstrap forecast-comparison inference that treats a pandemic as
one event rather than thousands of cells. The next break will not announce
itself either; the protocol is built to score it. The repair belongs to
the models that come after.


\section*{Reproducibility statement}
Code, the pre-registration document, the data manifest
(\texttt{data/MANIFEST.sha256}) and the \pkg{StMoMo} oracle dump are in the
project repository, whose public URL is released on acceptance. HMD data are not
redistributed; the manifest pins the vintage so any table can be tied to the
exact input files. Seeds are recorded per grid cell and per ensemble member.

\appendix
\section{Secondary tables}
\label{app:secondary}

\subsection{Proper scores for conformal cells}
Conformal wrappers emit intervals, not distributions; the harness draws
uniformly inside each interval so that the single sample interface is
preserved. CRPS, log score and PIT computed from those draws therefore
describe a placeholder density. They are carried per cell in the released
results files under a \texttt{secondary} flag, are read only against other
conformal cells, and are never ranked against distributional mechanisms
(Section~\ref{sec:design-mechanisms}). They are deliberately not printed
beside the distributional scores of Section~\ref{sec:results}: a reader who
compares a conformal cell's CRPS with a native cell's CRPS is comparing a
convention with a forecast, and a shared table invites exactly that. The
same convention governs the derived quantities, whose $e_0$, $e_{65}$ and
annuity quantiles for a conformal cell come from the filler draws and are
reported as not applicable throughout (Addendum~2, \S3).

\subsection{Robustness}
The pre-registered secondary runs re-score the headline coverage contrasts
on pre-specified subsets of the same fits rather than refitting: the shift
regime with 2024 dropped (provisional-revision sensitivity), with USA and
CHL dropped, and split into register-based versus census-based populations;
the placebo regime without GBR\_SCO, on the neutral stratum only, and with
Denmark's 1921--1922 removed (Addendum~1). Each slice is recomputed by
\texttt{scripts/sensitivities.py} into \texttt{results/sensitivities.json},
released with the code: it carries, for every family--mechanism arm in every
slice, mean 95\% coverage over the rows the slice retains, joint path
coverage wherever the slice keeps the whole path, and the row and
error-class counts behind each mean. Two limits are worth stating plainly.
A slice that drops a horizon cannot report joint path coverage, because the
runner emits one indicator for the whole path rather than the per-horizon
indicators a shortened path would need; and the registered age-cap
sensitivity, ages capped at 90 rather than 99, is not among the slices---the
age-band decomposition of Section~\ref{sec:results-age}, which scores
$0$--$24$, $25$--$64$ and $65$--$99$ separately, is what this paper reports
about the old-age stratum instead.

\bibliographystyle{plainnat}
\bibliography{references}

\section{Full ledger of cells that produced no valid row}
\label{app:infeasible}

Table~\ref{tab:infeasible} in Section~\ref{sec:results-infeasible}
summarises, by population, family and mechanism, the primary-grid cells
that produced no valid row. The ledger below is the unabridged version:
one row per affected cell, carrying the QA gate's class (structural,
method, machine), the full recorded reason string rather than an excerpt,
and the number of rows affected. It is generated by
\texttt{scripts/make\_tables.py} from the same QA output as the summary
table and is typeset as a \texttt{longtable} so that it may run over
several pages; nothing in it is typed by hand. The regime, population
count and total number of error rows are stated in its note.

\inputtable{tab-infeasible-full}


\end{document}
